\chapter{Relative Mealy Program Categories and Contextual Dynamic Proarrow Logic}
\label{ch:relative-mealy-program-categories}

\section{Overview}
\label{sec:relative-mealy-overview}

This chapter develops dynamic proarrow logic from the Eilenberg--Moore
semantics of internal categories regarded as monads in the bicategory of spans.
We use the span conventions, composition order, internal-category conventions,
Mealy morphisms, Mealy transformations, and representable-carrier conventions
fixed in Chapter~\ref{ch:spans-internal-categories-mealy}.

The central point of the chapter is that a Mealy morphism
\[
\Phi:\mathbb A\rightsquigarrow\mathbb B
\]
does not merely transform input arrows of \(\mathbb A\) into output arrows of
\(\mathbb B\). Since internal categories are monads in
\(\mathbf{Span}(\mathcal E)\), it also induces a stateful restriction operation
between their categories of right actions:
\[
\Phi^\ast:\mathbf{Act}(\mathbb B)\longrightarrow \mathbf{Act}(\mathbb A).
\]
Thus every right \(\mathbb B\)-action \(Y\) gives rise to a right
\(\mathbb A\)-action \(\Phi^\ast Y\).

The main object of this chapter is the relative program category
\[
\mathsf{Prog}_{\Phi,Y}:=\int_{\mathbb A}\Phi^\ast Y,
\]
the action category of the restricted action. This is the internal category of
executions of the Mealy adapter \(\Phi\) coupled to the downstream
\(\mathbb B\)-system \(Y\).

The terminal downstream action recovers the trace case. If
\[
\mathbf 1_{\mathbb B}
=
(B_0\xrightarrow{1_{B_0}}B_0)
\]
is the terminal right \(\mathbb B\)-action, then
\[
\mathsf{Prog}_{\Phi,\mathbf 1_{\mathbb B}}
\]
is the trace program category of the Mealy morphism itself. In this terminal
case, emitted \(\mathbb B\)-arrows are observed as labels. In the general
relative case, emitted \(\mathbb B\)-arrows are consumed by an actual
downstream \(\mathbb B\)-action.

The guiding slogan is:
\[
\boxed{
\text{relative Mealy dynamic proarrow logic is dynamic logic on the action
category of a \(B\)-behaviour restricted along a stateful Mealy morphism.}
}
\]

\section{Eilenberg--Moore actions of an internal category}
\label{sec:em-actions-internal-category}

Throughout this section, let
\[
\mathbb A=(A_0,A_1,s_A,t_A,e_A,m_A)
\]
be an internal category in \(\mathcal E\). By the conventions of
Chapter~\ref{ch:spans-internal-categories-mealy}, this is equivalently a
monad in \(\mathbf{Span}(\mathcal E)\) on the object \(A_0\), whose underlying
endospan is
\[
A_0\xleftarrow{s_A}A_1\xrightarrow{t_A}A_0.
\]

\begin{definition}[Right internal action]
\label{def:right-internal-action-relative}
A \textbf{right action} of \(\mathbb A\) consists of a morphism
\[
p:X\to A_0
\]
together with a map
\[
\alpha:
A_1\times_{t_A,A_0,p}X\to X
\]
such that
\[
p\alpha=s_A\pi_1.
\]
In generalized-element notation, if
\[
a:u\to v,
\qquad
p(x)=v,
\]
then
\[
\alpha(a,x)
\]
lies over \(u\).

The action is required to satisfy the unit law
\[
\alpha(e_A(px),x)=x
\]
and the multiplication law
\[
\alpha(m_A(a,b),x)=\alpha(a,\alpha(b,x)),
\]
for
\[
a:u\to v,
\qquad
b:v\to w,
\qquad
p(x)=w.
\]
Here \(m_A(a,b):u\to w\) uses the composition order fixed in
Chapter~\ref{ch:spans-internal-categories-mealy}.
\end{definition}

\begin{definition}[Equivariant map]
\label{def:equivariant-map-relative}
Let
\[
(X\xrightarrow{p}A_0,\alpha)
\qquad\text{and}\qquad
(Y\xrightarrow{q}A_0,\beta)
\]
be right \(\mathbb A\)-actions. An \(\mathbb A\)-equivariant map
\[
f:(X,p,\alpha)\to(Y,q,\beta)
\]
is a morphism
\[
f:X\to Y
\]
such that
\[
qf=p
\]
and
\[
f(\alpha(a,x))=\beta(a,f(x)).
\]
\end{definition}

\begin{definition}[Category of actions]
\label{def:act-category-relative}
The \textbf{category of right actions} of \(\mathbb A\), denoted
\[
\mathbf{Act}(\mathbb A),
\]
is the external category whose objects are right internal \(\mathbb A\)-actions
and whose morphisms are equivariant maps.
\end{definition}

\begin{proposition}[Actions as Eilenberg--Moore algebras]
\label{prop:actions-are-em-algebras}
The internal category \(\mathbb A\), regarded as a monad in
\(\mathbf{Span}(\mathcal E)\), induces an ordinary monad
\[
T_{\mathbb A}:\mathcal E/A_0\to\mathcal E/A_0
\]
given by
\[
T_{\mathbb A}(X\xrightarrow{p}A_0)
=
\left(
A_1\times_{t_A,A_0,p}X
\xrightarrow{s_A\pi_1}
A_0
\right).
\]
The category \(\mathbf{Act}(\mathbb A)\) is the ordinary
Eilenberg--Moore category of this monad. Equivalently, it is the external
category of right \(\mathbb A\)-modules for the span monad.
\end{proposition}

\begin{proof}
A \(T_{\mathbb A}\)-algebra structure on
\[
p:X\to A_0
\]
is precisely a map
\[
\alpha:
A_1\times_{t_A,A_0,p}X\to X
\]
over \(A_0\), where the output lies over the source of the acting arrow:
\[
p\alpha=s_A\pi_1.
\]
The unit and multiplication axioms for the monad \(T_{\mathbb A}\) are exactly
the unit and multiplication laws in
Definition~\ref{def:right-internal-action-relative}. Morphisms of
Eilenberg--Moore algebras are precisely the equivariant maps of
Definition~\ref{def:equivariant-map-relative}.
\end{proof}

\begin{remark}
\label{rem:external-em-versus-internal-action-category}
The category \(\mathbf{Act}(\mathbb A)\) is an external category of internal
actions. It should not be confused with the internal action category associated
to a particular action, which will be denoted
\[
\int_{\mathbb A}X.
\]
\end{remark}

\section{Action categories and input-discrete functors}
\label{sec:action-categories-input-discrete}

Let
\[
(X\xrightarrow{p}A_0,\alpha)
\]
be a right \(\mathbb A\)-action.

\begin{definition}[Action category]
\label{def:action-category-relative}
The \textbf{action category} of \(X\), denoted
\[
\int_{\mathbb A}X,
\]
is the internal category defined as follows:
\[
\left(\int_{\mathbb A}X\right)_0=X,
\]
\[
\left(\int_{\mathbb A}X\right)_1
=
A_1\times_{t_A,A_0,p}X.
\]
The source and target maps are
\[
s(a,x)=x,
\qquad
t(a,x)=\alpha(a,x).
\]
The identity map is
\[
e(x)=(e_A(px),x).
\]
Composition is given by
\[
\bigl((b,x),(a,\alpha(b,x))\bigr)
\longmapsto
(m_A(a,b),x).
\]
Thus \(m_A(a,b)\) denotes the composite of
\[
a:u\to v,
\qquad
b:v\to w,
\]
while the right action executes \(b\) at \(x\) and then executes \(a\) at the
updated state.
\end{definition}

\begin{theorem}[Action categories]
\label{thm:action-category-is-internal-category}
For every right \(\mathbb A\)-action
\[
(X\xrightarrow{p}A_0,\alpha),
\]
the data of Definition~\ref{def:action-category-relative} define an internal
category
\[
\int_{\mathbb A}X.
\]
Equivalently, the span
\[
X
\xleftarrow{\pi_2}
A_1\times_{t_A,A_0,p}X
\xrightarrow{\alpha}
X
\]
is a monoid object in
\[
\operatorname{EndSpan}_{\mathcal E}(X).
\]
\end{theorem}

\begin{proof}
The source and target maps are well-defined by the typing equation
\[
p\alpha=s_A\pi_1.
\]
The identity map is well-defined because
\[
t_A(e_A(px))=px.
\]
For composition, a composable pair has the form
\[
\bigl((b,x),(a,\alpha(b,x))\bigr),
\]
where
\[
a:u\to v,
\qquad
b:v\to w,
\qquad
p(x)=w.
\]
The composite is
\[
(m_A(a,b),x).
\]
This is an arrow of the action category because
\[
t_A(m_A(a,b))=t_A(b)=p(x).
\]
Its target is
\[
\alpha(m_A(a,b),x),
\]
which equals
\[
\alpha(a,\alpha(b,x))
\]
by the action multiplication law. Thus the composite has the expected source
and target. The unit laws follow from the unit laws of \(\mathbb A\) and the
action unit law. Associativity follows from associativity of \(m_A\) and the
action multiplication law.

The equivalent endospan-monoid statement is the standard equivalence between
internal categories with object object \(X\) and monoid objects in
\(\operatorname{EndSpan}_{\mathcal E}(X)\).
\end{proof}

\begin{remark}
\label{rem:action-category-execution}
The category \(\int_{\mathbb A}X\) is the internal category of executions of
the Eilenberg--Moore algebra \(X\). An arrow is an admissible pair
\[
(a,x),
\]
where \(a\) is an arrow of \(\mathbb A\) whose target matches the state \(x\).
Executing \(a\) at \(x\) changes the state to \(\alpha(a,x)\).
\end{remark}

The action category has a canonical internal functor
\[
\pi_X:\int_{\mathbb A}X\to\mathbb A^{\mathrm{op}}.
\]
On objects it is
\[
p:X\to A_0,
\]
and on arrows it is the projection
\[
\pi_1:A_1\times_{t_A,A_0,p}X\to A_1.
\]

\begin{definition}[Input-discrete functor]
\label{def:input-discrete-relative}
Let
\[
J:\mathbb X\to\mathbb A^{\mathrm{op}}
\]
be an internal functor with object map
\[
J_0:X_0\to A_0.
\]
We say that \(J\) is \textbf{input-discrete} if the square
\[
\begin{tikzcd}
X_1 \arrow[r,"J_1"] \arrow[d,"s_X"']
&
A_1 \arrow[d,"t_A"]
\\
X_0 \arrow[r,"J_0"']
&
A_0
\end{tikzcd}
\]
is a pullback.
\end{definition}

\begin{theorem}[Actions as input-discrete functors]
\label{thm:actions-as-input-discrete-relative}
To give a right \(\mathbb A\)-action
\[
(X\xrightarrow{p}A_0,\alpha)
\]
is equivalently to give an internal category \(\mathbb X\) with object object
\(X\), together with an input-discrete internal functor
\[
J:\mathbb X\to\mathbb A^{\mathrm{op}}
\]
whose object map is \(p\).
\end{theorem}

\begin{proof}
Starting from a right action, the action category
\[
\int_{\mathbb A}X
\]
has arrow object
\[
A_1\times_{t_A,A_0,p}X.
\]
The canonical functor
\[
\pi_X:\int_{\mathbb A}X\to\mathbb A^{\mathrm{op}}
\]
has arrow map \(\pi_1\). The square in
Definition~\ref{def:input-discrete-relative} is exactly the pullback square
defining \(A_1\times_{t_A,A_0,p}X\). Hence \(\pi_X\) is input-discrete.

Conversely, suppose
\[
J:\mathbb X\to\mathbb A^{\mathrm{op}}
\]
is input-discrete with object map \(p:X\to A_0\). Since the square in
Definition~\ref{def:input-discrete-relative} is a pullback, there is a
canonical isomorphism
\[
X_1\cong A_1\times_{t_A,A_0,p}X.
\]
Transporting the target map
\[
t_X:X_1\to X
\]
across this isomorphism gives a map
\[
\alpha:A_1\times_{t_A,A_0,p}X\to X.
\]
The target compatibility of the functor
\[
J:\mathbb X\to\mathbb A^{\mathrm{op}}
\]
says precisely
\[
p\alpha=s_A\pi_1.
\]
Identity preservation gives the action unit law, and composition preservation
gives the action multiplication law. The two constructions are inverse up to
the canonical pullback isomorphism.
\end{proof}

\section{Mealy restriction of Eilenberg--Moore actions}
\label{sec:mealy-restriction-em-actions}

Let
\[
\Phi=(P,\delta,\omega):\mathbb A\rightsquigarrow\mathbb B
\]
be a Mealy morphism as in Chapter~\ref{ch:spans-internal-categories-mealy},
with carrier span
\[
A_0\xleftarrow{p}P\xrightarrow{q}B_0.
\]
Thus \(\delta\) is the state-update component and \(\omega\) is the output
component.

Let
\[
Y=(Y\xrightarrow{r}B_0,\sigma)
\]
be a right \(\mathbb B\)-action. Thus
\[
\sigma:B_1\times_{t_B,B_0,r}Y\to Y
\]
satisfies the right-action laws for \(\mathbb B\).

\begin{definition}[Restriction of a downstream action]
\label{def:mealy-restriction-action}
Given \(\Phi\) and \(Y\), define
\[
\Phi^\ast Y
\]
to be the object
\[
P\times_{q,B_0,r}Y
\]
over \(A_0\) by
\[
P\times_{B_0}Y
\xrightarrow{\pi_P}
P
\xrightarrow{p}
A_0.
\]
The action map is
\[
A_1\times_{t_A,A_0,p\pi_P}(P\times_{B_0}Y)
\to
P\times_{B_0}Y
\]
given by
\[
a\cdot(x,y)
=
\left(
\delta(a,x),
\sigma(\omega(a,x),y)
\right).
\]
\end{definition}

\begin{theorem}[Mealy restriction of actions]
\label{thm:mealy-restriction-functor}
Every Mealy morphism
\[
\Phi:\mathbb A\rightsquigarrow\mathbb B
\]
induces a functor
\[
\Phi^\ast:\mathbf{Act}(\mathbb B)\to\mathbf{Act}(\mathbb A).
\]
\end{theorem}

\begin{proof}
Let
\[
Y=(Y\xrightarrow{r}B_0,\sigma)
\]
be a right \(\mathbb B\)-action. We first check that
\(\Phi^\ast Y\) is a right \(\mathbb A\)-action.

The displayed action is well-defined because the Mealy typing equation gives
\[
t_B\omega(a,x)=q(x)=r(y),
\]
so \(\sigma(\omega(a,x),y)\) is defined. Moreover,
\[
r\sigma(\omega(a,x),y)
=
s_B\omega(a,x)
=
q\delta(a,x).
\]
Thus
\[
\left(
\delta(a,x),
\sigma(\omega(a,x),y)
\right)
\]
lies in \(P\times_{B_0}Y\).

The projection to \(A_0\) sends this pair to
\[
p\delta(a,x)=s_A(a),
\]
so the action has the correct typing.

The unit law follows from the Mealy unit laws and the unit law of the
\(\mathbb B\)-action:
\[
e_A(px)\cdot(x,y)
=
\left(
\delta(e_A(px),x),
\sigma(\omega(e_A(px),x),y)
\right)
=
(x,y).
\]
The multiplication law follows from the Mealy multiplication laws and the
multiplication law of the \(\mathbb B\)-action:
\[
m_A(a,b)\cdot(x,y)
=
a\cdot\bigl(b\cdot(x,y)\bigr).
\]
Thus \(\Phi^\ast Y\) is a right \(\mathbb A\)-action.

If
\[
f:Y\to Z
\]
is an equivariant map of right \(\mathbb B\)-actions, define
\[
\Phi^\ast f:P\times_{B_0}Y\to P\times_{B_0}Z
\]
by
\[
\Phi^\ast f(x,y)=(x,f(y)).
\]
This is well-defined because \(f\) lies over \(B_0\). It is equivariant by
equivariance of \(f\):
\[
f(\sigma_Y(\omega(a,x),y))
=
\sigma_Z(\omega(a,x),f(y)).
\]
Functoriality is immediate.
\end{proof}

\begin{remark}
\label{rem:mealy-restriction-slogan}
A Mealy morphism is a stateful restriction-of-scalars operation on actions. A
\(\mathbb B\)-system becomes an \(\mathbb A\)-system by placing the Mealy
adapter in front of it:
\[
\mathbb A\text{-input}
\longrightarrow
\Phi
\longrightarrow
\mathbb B\text{-output}
\longrightarrow
Y.
\]
\end{remark}

\section{Relative Mealy program categories}
\label{sec:relative-mealy-program-categories}

Fix a Mealy morphism
\[
\Phi=(P,\delta,\omega):\mathbb A\rightsquigarrow\mathbb B
\]
and a right \(\mathbb B\)-action
\[
Y=(Y\xrightarrow{r}B_0,\sigma).
\]

\begin{definition}[Relative state object]
\label{def:relative-state-object}
The \textbf{relative state object} of \(\Phi\) at \(Y\) is
\[
S_{\Phi,Y}:=P\times_{q,B_0,r}Y.
\]
A generalized element is written
\[
(x,y)
\]
with
\[
q(x)=r(y).
\]
\end{definition}

\begin{definition}[Relative primitive execution object]
\label{def:relative-primitive-execution-object}
The \textbf{relative primitive execution object} is
\[
M_{\Phi,Y}
:=
A_1\times_{t_A,A_0,p\pi_P}S_{\Phi,Y}.
\]
A generalized element is a triple
\[
(a,x,y)
\]
such that
\[
t_A(a)=p(x),
\qquad
q(x)=r(y).
\]
\end{definition}

\begin{remark}
\label{rem:primitive-relative-to-input-category}
The adjective ``primitive'' is relative to the chosen internal category
\(\mathbb A\). If arrows of \(\mathbb A\) already represent composite
programs, then primitive executions here consume such composite arrows as
single steps. Generator-level primitive steps require restricting
\(M_{\Phi,Y}\) to a chosen subobject, for instance to a subobject of
\(A_1\) pulled back along the input projection.
\end{remark}

\begin{definition}[Relative Mealy program category]
\label{def:relative-mealy-program-category}
The \textbf{relative Mealy program category} of \(\Phi\) at \(Y\) is
\[
\mathsf{Prog}_{\Phi,Y}
:=
\int_{\mathbb A}\Phi^\ast Y.
\]
Explicitly,
\[
(\mathsf{Prog}_{\Phi,Y})_0=S_{\Phi,Y},
\]
\[
(\mathsf{Prog}_{\Phi,Y})_1=M_{\Phi,Y}.
\]
The source and target maps are
\[
s_{\Phi,Y}(a,x,y)=(x,y),
\]
and
\[
t_{\Phi,Y}(a,x,y)
=
\left(
\delta(a,x),
\sigma(\omega(a,x),y)
\right).
\]
The identity is
\[
e_{\Phi,Y}(x,y)=(e_A(px),x,y).
\]
Composition is
\[
\bigl((b,x,y),(a,\delta(b,x),\sigma(\omega(b,x),y))\bigr)
\longmapsto
(m_A(a,b),x,y).
\]
\end{definition}

\begin{theorem}[Relative program categories]
\label{thm:relative-program-category}
For every Mealy morphism
\[
\Phi:\mathbb A\rightsquigarrow\mathbb B
\]
and every right \(\mathbb B\)-action \(Y\), the data of
Definition~\ref{def:relative-mealy-program-category} define an internal
category
\[
\mathsf{Prog}_{\Phi,Y}.
\]
Equivalently, the span
\[
\mathbf m_{\Phi,Y}
=
\left(
S_{\Phi,Y}
\xleftarrow{s_{\Phi,Y}}
M_{\Phi,Y}
\xrightarrow{t_{\Phi,Y}}
S_{\Phi,Y}
\right)
\]
is a monoid object in
\[
\operatorname{EndSpan}_{\mathcal E}(S_{\Phi,Y}).
\]
\end{theorem}

\begin{proof}
By Theorem~\ref{thm:mealy-restriction-functor}, \(\Phi^\ast Y\) is a right
\(\mathbb A\)-action. Therefore its action category
\[
\int_{\mathbb A}\Phi^\ast Y
\]
is an internal category by
Theorem~\ref{thm:action-category-is-internal-category}. The displayed formulas
are the explicit instance of the general action-category construction.
\end{proof}

\begin{remark}
\label{rem:relative-program-category-interpretation}
The category \(\mathsf{Prog}_{\Phi,Y}\) is the execution category of the
Mealy adapter \(\Phi\) coupled to the downstream \(\mathbb B\)-system \(Y\).
A state is a pair
\[
(x,y)\in P\times_{B_0}Y.
\]
An \(A\)-input updates \(x\) by \(\delta\), emits a \(B\)-arrow by \(\omega\),
and then lets \(Y\) consume that \(B\)-arrow through its action \(\sigma\).
\end{remark}

\section{The terminal trace case}
\label{sec:terminal-trace-case}

The relative theory contains the trace semantics of a Mealy morphism as the
special case obtained by taking the terminal downstream action.

\begin{definition}[Terminal right action]
\label{def:terminal-right-action}
Let \(\mathbb B\) be an internal category. Its \textbf{terminal right action}
is
\[
\mathbf 1_{\mathbb B}
=
(B_0\xrightarrow{1_{B_0}}B_0,\tau_B),
\]
where
\[
\tau_B:
B_1\times_{t_B,B_0,1_{B_0}}B_0\to B_0
\]
is given by
\[
\tau_B(b,t_Bb)=s_Bb.
\]
\end{definition}

\begin{proposition}
\label{prop:terminal-action-category}
The action category of the terminal right action is canonically isomorphic to
the opposite internal category:
\[
\int_{\mathbb B}\mathbf 1_{\mathbb B}
\cong
\mathbb B^{\mathrm{op}}.
\]
\end{proposition}

\begin{proof}
The object object is \(B_0\). The arrow object is
\[
B_1\times_{t_B,B_0,1_{B_0}}B_0\cong B_1.
\]
Under this identification, the source of the action category is \(t_B\), and
the target is \(s_B\). Thus the source and target are those of
\(\mathbb B^{\mathrm{op}}\). The identity and composition maps are also those
of \(\mathbb B^{\mathrm{op}}\).
\end{proof}

\begin{corollary}[Terminal trace program category]
\label{cor:terminal-trace-program-category}
Let
\[
\Phi=(P,\delta,\omega):\mathbb A\rightsquigarrow\mathbb B
\]
be a Mealy morphism. Then
\[
\Phi^\ast(\mathbf 1_{\mathbb B})
\cong
(P\xrightarrow{p}A_0,\delta)
\]
as a right \(\mathbb A\)-action. Consequently
\[
\mathsf{Prog}_{\Phi,\mathbf 1_{\mathbb B}}
\]
has object object \(P\), arrow object
\[
A_1\times_{t_A,A_0,p}P,
\]
source and target
\[
s(a,x)=x,
\qquad
t(a,x)=\delta(a,x).
\]
\end{corollary}

\begin{proof}
Since the underlying object of \(\mathbf 1_{\mathbb B}\) is \(B_0\) over
itself, there is a canonical isomorphism
\[
P\times_{q,B_0,1_{B_0}}B_0\cong P.
\]
Under this isomorphism, the induced action of \(\mathbb A\) is
\[
a\cdot x=\delta(a,x).
\]
The relative program category therefore specializes to the displayed trace
program category.
\end{proof}

\begin{definition}[Trace program notation]
\label{def:trace-program-notation}
We write
\[
\mathsf{Prog}_\Phi
:=
\mathsf{Prog}_{\Phi,\mathbf 1_{\mathbb B}},
\]
and
\[
\mathbf m_\Phi
:=
\mathbf m_{\Phi,\mathbf 1_{\mathbb B}}
\]
for the terminal trace program category and its primitive execution span.
\end{definition}

\begin{remark}
\label{rem:terminal-case-trace-interpretation}
The terminal case is trace semantics. The emitted \(B\)-arrow is recorded as
an output label, but it is not consumed by a nontrivial downstream system.
The general relative case is contextual semantics: the emitted \(B\)-arrow is
executed in a chosen \(B\)-action \(Y\).
\end{remark}

\section{Input projection, downstream comparison, and combined labelling}
\label{sec:input-downstream-combined-labelling}

Let
\[
\Phi:\mathbb A\rightsquigarrow\mathbb B
\]
be a Mealy morphism and let \(Y\) be a right \(\mathbb B\)-action.

\begin{definition}[Input projection]
\label{def:relative-input-projection}
The \textbf{input projection}
\[
I_{\Phi,Y}:\mathsf{Prog}_{\Phi,Y}\to\mathbb A^{\mathrm{op}}
\]
is defined on objects by
\[
I_{\Phi,Y}(x,y)=p(x)
\]
and on arrows by
\[
I_{\Phi,Y}(a,x,y)=a.
\]
\end{definition}

\begin{definition}[Downstream comparison functor]
\label{def:downstream-comparison-functor}
The \textbf{downstream comparison functor}
\[
\Lambda_{\Phi,Y}:\mathsf{Prog}_{\Phi,Y}\to\int_{\mathbb B}Y
\]
is defined on objects by
\[
\Lambda_{\Phi,Y}(x,y)=y
\]
and on arrows by
\[
\Lambda_{\Phi,Y}(a,x,y)=(\omega(a,x),y).
\]
\end{definition}

\begin{theorem}
\label{thm:input-and-downstream-functors}
The map \(I_{\Phi,Y}\) is an input-discrete internal functor
\[
\mathsf{Prog}_{\Phi,Y}\to\mathbb A^{\mathrm{op}}.
\]
The map \(\Lambda_{\Phi,Y}\) is an internal functor
\[
\mathsf{Prog}_{\Phi,Y}\to\int_{\mathbb B}Y.
\]
\end{theorem}

\begin{proof}
The functor \(I_{\Phi,Y}\) is the canonical input projection associated to the
right \(\mathbb A\)-action \(\Phi^\ast Y\). It is input-discrete by
Theorem~\ref{thm:actions-as-input-discrete-relative}.

For \(\Lambda_{\Phi,Y}\), first note that the arrow map is well-defined:
\[
t_B\omega(a,x)=q(x)=r(y).
\]
Thus
\[
(\omega(a,x),y)
\]
is an arrow of the action category \(\int_{\mathbb B}Y\).

The source is preserved because
\[
s(\omega(a,x),y)=y.
\]
The target is preserved because
\[
t(\omega(a,x),y)
=
\sigma(\omega(a,x),y),
\]
which is the \(Y\)-component of \(t_{\Phi,Y}(a,x,y)\).

Identity preservation follows from the Mealy unit law for \(\omega\).

For composition, consider a composable pair
\[
\bigl((b,x,y),(a,\delta(b,x),\sigma(\omega(b,x),y))\bigr).
\]
Its composite in \(\mathsf{Prog}_{\Phi,Y}\) is
\[
(m_A(a,b),x,y).
\]
Applying \(\Lambda_{\Phi,Y}\) gives
\[
(\omega(m_A(a,b),x),y).
\]
On the other hand, applying \(\Lambda_{\Phi,Y}\) first gives the composable
pair
\[
\bigl(
(\omega(b,x),y),
(\omega(a,\delta(b,x)),\sigma(\omega(b,x),y))
\bigr)
\]
in \(\int_{\mathbb B}Y\). Its composite is
\[
\bigl(
m_B(\omega(a,\delta(b,x)),\omega(b,x)),y
\bigr).
\]
These agree by the Mealy output multiplication law
\[
\omega(m_A(a,b),x)
=
m_B(\omega(a,\delta(b,x)),\omega(b,x)).
\]
Thus \(\Lambda_{\Phi,Y}\) preserves composition.
\end{proof}

\begin{corollary}[Combined relative labelling]
\label{cor:combined-relative-labelling}
If the product internal category
\[
\mathbb A^{\mathrm{op}}\times \int_{\mathbb B}Y
\]
exists, then \(I_{\Phi,Y}\) and \(\Lambda_{\Phi,Y}\) assemble into a single
internal functor
\[
L_{\Phi,Y}:
\mathsf{Prog}_{\Phi,Y}
\longrightarrow
\mathbb A^{\mathrm{op}}\times \int_{\mathbb B}Y.
\]
On objects,
\[
L_{\Phi,Y}(x,y)=(p x,y),
\]
and on arrows,
\[
L_{\Phi,Y}(a,x,y)=\bigl(a,(\omega(a,x),y)\bigr).
\]
\end{corollary}

\begin{proof}
This is the universal property of the product internal category.
\end{proof}

\begin{corollary}[Terminal output labelling]
\label{cor:terminal-output-labelling}
For \(Y=\mathbf 1_{\mathbb B}\), the downstream comparison functor becomes the
output labelling functor
\[
\Omega_\Phi:
\mathsf{Prog}_\Phi\to\mathbb B^{\mathrm{op}}.
\]
It is given on objects by
\[
x\mapsto q(x)
\]
and on arrows by
\[
(a,x)\mapsto \omega(a,x).
\]
If products exist, the combined labelling is
\[
L_\Phi:
\mathsf{Prog}_\Phi
\to
\mathbb A^{\mathrm{op}}\times \mathbb B^{\mathrm{op}}.
\]
\end{corollary}

\begin{proof}
By Proposition~\ref{prop:terminal-action-category},
\[
\int_{\mathbb B}\mathbf 1_{\mathbb B}
\cong
\mathbb B^{\mathrm{op}}.
\]
Under this identification, \(\Lambda_{\Phi,\mathbf 1_{\mathbb B}}\) is exactly
the displayed output functor.
\end{proof}

\section{Functoriality in the downstream behaviour}
\label{sec:functoriality-downstream-behaviour}

The assignment
\[
Y\longmapsto \mathsf{Prog}_{\Phi,Y}
\]
is functorial in the downstream \(\mathbb B\)-action.

\begin{definition}[Induced functor on relative program categories]
\label{def:induced-functor-program-categories}
Let
\[
f:Y\to Z
\]
be a morphism of right \(\mathbb B\)-actions. Define
\[
\mathsf{Prog}_{\Phi,f}:
\mathsf{Prog}_{\Phi,Y}\to\mathsf{Prog}_{\Phi,Z}
\]
on objects by
\[
(x,y)\mapsto(x,f(y))
\]
and on arrows by
\[
(a,x,y)\mapsto(a,x,f(y)).
\]
\end{definition}

\begin{theorem}
\label{thm:program-categories-functorial-in-Y}
For each Mealy morphism
\[
\Phi:\mathbb A\rightsquigarrow\mathbb B,
\]
the assignment
\[
Y\longmapsto \mathsf{Prog}_{\Phi,Y}
\]
extends to a functor
\[
\mathbf{Act}(\mathbb B)
\longrightarrow
\mathbf{Cat}_{\mathrm{int}}(\mathcal E).
\]
More precisely, it factors as
\[
\mathbf{Act}(\mathbb B)
\xrightarrow{\Phi^\ast}
\mathbf{Act}(\mathbb A)
\xrightarrow{\int_{\mathbb A}}
\mathbf{Cat}_{\mathrm{int}}(\mathcal E)/\mathbb A^{\mathrm{op}}.
\]
\end{theorem}

\begin{proof}
The map on objects and arrows is well-defined because \(f\) lies over \(B_0\).
It preserves sources immediately. It preserves targets because \(f\) is
\(\mathbb B\)-equivariant:
\[
f(\sigma_Y(\omega(a,x),y))
=
\sigma_Z(\omega(a,x),f(y)).
\]
It preserves identities and composition by the displayed formulas. Functoriality
in \(f\) is immediate. The factorization is the definition of
\(\mathsf{Prog}_{\Phi,Y}\) as the action category of \(\Phi^\ast Y\).
\end{proof}

\begin{remark}
\label{rem:change-of-downstream-model}
A morphism
\[
f:Y\to Z
\]
of downstream actions is a change of downstream model. It may be read as a
refinement, abstraction, observation, or quotient of the receiving
\(\mathbb B\)-system. The induced functor
\[
\mathsf{Prog}_{\Phi,Y}\to\mathsf{Prog}_{\Phi,Z}
\]
transports relative executions along that change of model.
\end{remark}

\section{Mealy transformations and relative labelled functors}
\label{sec:mealy-transformations-relative-labelled-functors}

Let
\[
\Phi=(P,\delta_P,\omega_P),
\qquad
\Psi=(Q,\delta_Q,\omega_Q)
:
\mathbb A\rightsquigarrow\mathbb B
\]
be Mealy morphisms, and let
\[
\theta:\Phi\Rightarrow\Psi
\]
be a Mealy transformation in the sense of
Chapter~\ref{ch:spans-internal-categories-mealy}. Thus \(\theta:P\to Q\) lies
over both \(A_0\) and \(B_0\), commutes with the transition maps, and preserves
the output maps.

\begin{definition}[Relative functor induced by a Mealy transformation]
\label{def:relative-functor-induced-mealy-transformation}
For every right \(\mathbb B\)-action \(Y\), define
\[
\theta_Y:S_{\Phi,Y}\to S_{\Psi,Y}
\]
by
\[
\theta_Y(x,y)=(\theta x,y).
\]
On primitive executions, define
\[
\Theta_{Y,1}:M_{\Phi,Y}\to M_{\Psi,Y}
\]
by
\[
\Theta_{Y,1}(a,x,y)=(a,\theta x,y).
\]
\end{definition}

\begin{proposition}[Relative functoriality of Mealy transformations]
\label{prop:relative-functoriality-mealy-transformations}
The maps \(\theta_Y\) and \(\Theta_{Y,1}\) define an internal functor
\[
\Theta_Y:
\mathsf{Prog}_{\Phi,Y}\to\mathsf{Prog}_{\Psi,Y}.
\]
Moreover, \(\Theta_Y\) commutes with both the input projections and the
downstream comparison functors:
\[
I_{\Psi,Y}\Theta_Y=I_{\Phi,Y},
\qquad
\Lambda_{\Psi,Y}\Theta_Y=\Lambda_{\Phi,Y}.
\]
The assignment
\[
Y\longmapsto \Theta_Y
\]
is natural in \(Y\).
\end{proposition}

\begin{proof}
Since \(\theta\) lies over \(A_0\) and \(B_0\), the displayed maps are
well-defined. Source preservation is immediate. Target preservation follows
from the transition compatibility of the Mealy transformation:
\[
\theta(\delta_P(a,x))=\delta_Q(a,\theta x).
\]
The downstream \(Y\)-component is preserved because the output compatibility
gives
\[
\omega_P(a,x)=\omega_Q(a,\theta x).
\]
Thus
\[
\sigma(\omega_P(a,x),y)
=
\sigma(\omega_Q(a,\theta x),y).
\]
Identity and composition preservation follow from the displayed formulas and
the corresponding structure in the two relative program categories.

Commutation with input projections is immediate because the \(A_1\)-component
of an arrow is unchanged. Commutation with downstream comparison follows from
the output compatibility equation. Naturality in \(Y\) follows from the formula
\[
(x,y)\mapsto(\theta x,y).
\]
\end{proof}

\begin{corollary}[Terminal labelled functor theorem]
\label{cor:terminal-transformations-labelled-functors}
In the terminal case \(Y=\mathbf 1_{\mathbb B}\), a Mealy transformation
\[
\theta:\Phi\Rightarrow\Psi
\]
induces an internal functor
\[
\Theta:
\mathsf{Prog}_\Phi\to\mathsf{Prog}_\Psi
\]
commuting with both input and output labellings:
\[
I_\Psi\Theta=I_\Phi,
\qquad
\Omega_\Psi\Theta=\Omega_\Phi.
\]
Conversely, an internal functor
\[
\Theta:
\mathsf{Prog}_\Phi\to\mathsf{Prog}_\Psi
\]
commuting with both the input and output labellings determines a Mealy
transformation
\[
\Phi\Rightarrow\Psi.
\]
\end{corollary}

\begin{proof}
The forward direction is the terminal specialization of
Proposition~\ref{prop:relative-functoriality-mealy-transformations}.

Conversely, commutation on objects gives that the object map
\[
\theta:P\to Q
\]
lies over both \(A_0\) and \(B_0\), hence is a map of carrier spans.
Commutation with input labelling forces the arrow map to have the form
\[
(a,x)\mapsto(a,\theta x).
\]
Preservation of targets gives
\[
\theta(\delta_P(a,x))=\delta_Q(a,\theta x),
\]
and commutation with output labelling gives
\[
\omega_P(a,x)=\omega_Q(a,\theta x).
\]
These are precisely the Mealy transformation equations.
\end{proof}

\section{Labelled input-discrete characterizations}
\label{sec:labelled-input-discrete-characterizations}

The terminal trace case has a sharp universal characterization: Mealy
morphisms are labelled input-discrete categories. The relative construction
then interprets such labelled input-discrete categories in a downstream
\(\mathbb B\)-action.

\begin{definition}[Terminal labelled input-discrete category]
\label{def:terminal-labelled-input-discrete-category}
A \textbf{terminal labelled input-discrete category from \(\mathbb A\) to
\(\mathbb B\)} consists of:
\begin{enumerate}
\item an internal category \(\mathbb P\);
\item an input-discrete internal functor
\[
I:\mathbb P\to\mathbb A^{\mathrm{op}};
\]
\item an internal functor
\[
\Omega:\mathbb P\to\mathbb B^{\mathrm{op}}.
\]
\end{enumerate}
\end{definition}

\begin{theorem}[Mealy morphisms as labelled input-discrete categories]
\label{thm:mealy-as-labelled-input-discrete-relative}
To give a Mealy morphism
\[
\Phi:\mathbb A\rightsquigarrow\mathbb B
\]
is equivalently to give a terminal labelled input-discrete category from
\(\mathbb A\) to \(\mathbb B\).

Under this equivalence, the labelled input-discrete category associated to
\(\Phi\) is
\[
\left(
\mathsf{Prog}_\Phi,
I_\Phi,
\Omega_\Phi
\right).
\]
\end{theorem}

\begin{proof}
Starting from a Mealy morphism \(\Phi\), the terminal trace program category
\[
\mathsf{Prog}_\Phi
\]
has input projection
\[
I_\Phi:\mathsf{Prog}_\Phi\to\mathbb A^{\mathrm{op}}
\]
and output labelling
\[
\Omega_\Phi:\mathsf{Prog}_\Phi\to\mathbb B^{\mathrm{op}}.
\]
The input projection is input-discrete by
Theorem~\ref{thm:actions-as-input-discrete-relative}. Hence
\[
(\mathsf{Prog}_\Phi,I_\Phi,\Omega_\Phi)
\]
is a terminal labelled input-discrete category.

Conversely, suppose given
\[
I:\mathbb P\to\mathbb A^{\mathrm{op}},
\qquad
\Omega:\mathbb P\to\mathbb B^{\mathrm{op}},
\]
with \(I\) input-discrete. Let \(P\) be the object object of \(\mathbb P\).
The object maps of \(I\) and \(\Omega\) give
\[
p:P\to A_0,
\qquad
q:P\to B_0.
\]
Since \(I\) is input-discrete, the arrow object of \(\mathbb P\) is canonically
isomorphic to
\[
A_1\times_{t_A,A_0,p}P.
\]
Transporting the target map of \(\mathbb P\) across this pullback
identification gives
\[
\delta:A_1\times_{t_A,A_0,p}P\to P.
\]
The arrow map of \(\Omega\), transported across the same identification, gives
\[
\omega:A_1\times_{t_A,A_0,p}P\to B_1.
\]
The functoriality of \(I\) gives the typing and action equations for
\(\delta\). The functoriality of \(\Omega\) gives the source-target equations
for \(\omega\), the output unit law, and the output multiplication law. These
are exactly the component laws of a Mealy morphism.

The two constructions are inverse up to the canonical pullback isomorphism.
\end{proof}

\begin{remark}
\label{rem:labelled-input-discrete-not-ordinary-2-theory}
This is an equivalence of the data of Mealy morphisms with the data of
labelled input-discrete internal categories. It should not be read as an
identification of the resulting transformation theory with the ordinary
2-categorical theory of internal functors and internal natural
transformations.
\end{remark}

\begin{definition}[Relative labelled input-discrete category over \(Y\)]
\label{def:relative-labelled-input-discrete-category}
Fix a right \(\mathbb B\)-action \(Y\). A \textbf{relative labelled
input-discrete category over \(Y\)} consists of an internal category
\(\mathbb P\), an input-discrete functor
\[
I:\mathbb P\to\mathbb A^{\mathrm{op}},
\]
and an internal functor
\[
\Lambda:\mathbb P\to\int_{\mathbb B}Y.
\]
\end{definition}

\begin{proposition}[Relative interpretation of labelled input-discrete categories]
\label{prop:relative-labelled-input-discrete-interpretation}
For every Mealy morphism
\[
\Phi:\mathbb A\rightsquigarrow\mathbb B
\]
and every right \(\mathbb B\)-action \(Y\), the triple
\[
\left(
\mathsf{Prog}_{\Phi,Y},
I_{\Phi,Y},
\Lambda_{\Phi,Y}
\right)
\]
is a relative labelled input-discrete category over \(Y\).
\end{proposition}

\begin{proof}
This is exactly Theorem~\ref{thm:input-and-downstream-functors}.
\end{proof}

\begin{remark}
\label{rem:terminal-classification-relative-interpretation}
The terminal labelled input-discrete category classifies the Mealy morphism
itself. The relative labelled input-discrete category interprets that Mealy
morphism in a chosen downstream action \(Y\). Thus the terminal theory records
output labels, while the relative theory records the downstream transitions
caused by those outputs.
\end{remark}

\section{Ordinary functors as stateless Mealy morphisms}
\label{sec:ordinary-functors-stateless-mealy}

The relative theory includes ordinary internal functors as the stateless case.

Let
\[
F:\mathbb A\to\mathbb B
\]
be an internal functor, with object map
\[
F_0:A_0\to B_0
\]
and arrow map
\[
F_1:A_1\to B_1.
\]
As in Chapter~\ref{ch:spans-internal-categories-mealy}, \(F\) determines a
representable Mealy morphism whose carrier span is
\[
A_0\xleftarrow{1_{A_0}}A_0\xrightarrow{F_0}B_0.
\]

\begin{proposition}[Ordinary restriction of actions]
\label{prop:ordinary-functor-restriction-actions}
Let
\[
Y=(Y\xrightarrow{r}B_0,\sigma)
\]
be a right \(\mathbb B\)-action. For the stateless Mealy morphism induced by
\(F\), the restricted action is
\[
F^\ast Y
=
A_0\times_{F_0,B_0,r}Y
\]
with action
\[
a\cdot(v,y)
=
\bigl(s_A(a),\sigma(F_1(a),y)\bigr),
\]
where
\[
a:s_A(a)\to v
\]
and
\[
r(y)=F_0(v).
\]
\end{proposition}

\begin{proof}
This is the formula for Mealy restriction specialized to the representable
carrier
\[
A_0\xleftarrow{1}A_0\xrightarrow{F_0}B_0.
\]
The state component is forced to update from \(v=t_A(a)\) to \(s_A(a)\), and
the emitted \(B\)-arrow is \(F_1(a)\).
\end{proof}

\begin{corollary}[Terminal stateless case]
\label{cor:terminal-stateless-case}
For \(Y=\mathbf 1_{\mathbb B}\), one has
\[
\mathsf{Prog}_{F,\mathbf 1_{\mathbb B}}
\cong
\mathbb A^{\mathrm{op}}.
\]
Under this identification, the input projection is the identity functor on
\(\mathbb A^{\mathrm{op}}\), and the downstream comparison functor is
\[
F^{\mathrm{op}}:\mathbb A^{\mathrm{op}}\to\mathbb B^{\mathrm{op}}.
\]
\end{corollary}

\begin{proof}
In the terminal case,
\[
A_0\times_{F_0,B_0,1_{B_0}}B_0\cong A_0.
\]
The restricted action is the terminal \(\mathbb A\)-action. Its action category
is \(\mathbb A^{\mathrm{op}}\). The downstream comparison functor sends an
arrow \(a\) to \(F_1(a)\), viewed in the opposite category, hence it is
\(F^{\mathrm{op}}\).
\end{proof}

\begin{corollary}[Representable carriers as trivial input-discrete categories]
\label{cor:representable-trivial-input-discrete}
In the terminal trace case, a representable Mealy morphism induced by an
internal functor
\[
F:\mathbb A\to\mathbb B
\]
corresponds to the labelled input-discrete category for which the input
projection
\[
I_F:\mathsf{Prog}_F\to\mathbb A^{\mathrm{op}}
\]
is isomorphic to the identity functor on \(\mathbb A^{\mathrm{op}}\).
\end{corollary}

\begin{proof}
This follows from Corollary~\ref{cor:terminal-stateless-case}. Conversely, if
the input projection is isomorphic to the identity, then the carrier object is
isomorphic to \(A_0\) over \(A_0\), so the carrier span is representable and
the Mealy morphism is induced by an internal functor.
\end{proof}

\begin{remark}
\label{rem:natural-transformations-caveat-relative}
The stateless construction recovers internal functors at the level of objects
and 1-cells. The Mealy transformations between stateless carriers are generally
too strict to recover ordinary internal natural transformations. Thus the
identification with internal functors should be read at the level of objects
and morphisms, not as a recovery of the full ordinary 2-categorical theory of
internal categories, functors, and natural transformations.
\end{remark}

\section{Predicate transformers for spans}
\label{sec:relative-predicate-transformers-spans}

We now pass from the operational layer to the logical layer. The constructions
above require only pullbacks. To interpret modalities, we assume regular image
factorizations.

\begin{definition}[Regular category]
\label{def:regular-category-relative}
A category \(\mathcal E\) is \textbf{regular} if:
\begin{enumerate}
\item \(\mathcal E\) has finite limits;
\item every morphism factors as a regular epimorphism followed by a monomorphism;
\item regular epimorphisms are stable under pullback.
\end{enumerate}
\end{definition}

For the rest of this section, assume that \(\mathcal E\) is regular. For an
object \(X\), write
\[
\mathrm{Sub}(X)
\]
for the poset of subobjects of \(X\). For a morphism \(f:X\to Y\), write
\[
f^\ast:\mathrm{Sub}(Y)\to\mathrm{Sub}(X)
\]
for pullback along \(f\).

\begin{definition}[Existential image]
\label{def:existential-image-relative}
For a morphism
\[
f:X\to Y,
\]
define
\[
\exists_f:\mathrm{Sub}(X)\to\mathrm{Sub}(Y)
\]
by sending a subobject
\[
U\hookrightarrow X
\]
to the image of the composite
\[
U\to X\xrightarrow{f}Y.
\]
\end{definition}

\begin{proposition}
\label{prop:existential-left-adjoint-relative}
For every morphism \(f:X\to Y\) in a regular category, there is an adjunction
\[
\exists_f\dashv f^\ast.
\]
Thus, for \(U\in\mathrm{Sub}(X)\) and \(V\in\mathrm{Sub}(Y)\),
\[
\exists_f(U)\le V
\quad\Longleftrightarrow\quad
U\le f^\ast V.
\]
\end{proposition}

\begin{proof}
The subobject \(\exists_f(U)\) is the image of the composite \(U\to Y\). To
say that this image lies below \(V\hookrightarrow Y\) is to say that
\(U\to Y\) factors through \(V\). By the universal property of pullback, this
is equivalent to \(U\to X\) factoring through \(f^\ast V\).
\end{proof}

\begin{proposition}[Beck--Chevalley for regular images]
\label{prop:regular-beck-chevalley-relative}
For every pullback square
\[
\begin{tikzcd}
X' \arrow[r,"g'"] \arrow[d,"f'"']
&
X \arrow[d,"f"]
\\
Y' \arrow[r,"g"']
&
Y
\end{tikzcd}
\]
in a regular category, the Beck--Chevalley identity
\[
g^\ast\exists_f=\exists_{f'}(g')^\ast
\]
holds.
\end{proposition}

\begin{proof}
This is the stability of regular images under pullback. Pulling back the
image factorization of \(U\to X\xrightarrow{f}Y\) along \(g\) gives the image
factorization of the pulled-back subobject along \(f'\).
\end{proof}

\begin{definition}[Diamond modality of a span]
\label{def:span-diamond-relative}
Let
\[
R=
\left(
X\xleftarrow{\ell}R\xrightarrow{r}Y
\right)
\]
be a span in a regular category. Define
\[
\langle R\rangle:\mathrm{Sub}(Y)\to\mathrm{Sub}(X)
\]
by
\[
\langle R\rangle\varphi
:=
\exists_\ell(r^\ast\varphi).
\]
\end{definition}

\begin{proposition}[Functoriality of diamonds]
\label{prop:diamond-functoriality-relative}
For spans
\[
R:X\nrightarrow Y
\qquad\text{and}\qquad
S:Y\nrightarrow Z
\]
in a regular category,
\[
\langle 1_X\rangle=\mathrm{id}_{\mathrm{Sub}(X)}
\]
and
\[
\langle R;S\rangle
=
\langle R\rangle\circ\langle S\rangle.
\]
\end{proposition}

\begin{proof}
The identity law is immediate. For composition, write
\[
R=(X\xleftarrow{\ell}R\xrightarrow{r}Y),
\qquad
S=(Y\xleftarrow{m}S\xrightarrow{n}Z).
\]
The composite has apex \(R\times_Y S\), so
\[
\langle R;S\rangle\varphi
=
\exists_{\ell\pi_R}\bigl((n\pi_S)^\ast\varphi\bigr).
\]
Using functoriality of existential images and Beck--Chevalley for the
pullback square defining \(R\times_Y S\), this becomes
\[
\exists_\ell r^\ast\exists_m n^\ast\varphi
=
\langle R\rangle(\langle S\rangle\varphi).
\]
\end{proof}

For a morphism \(f:X\to Y\), define the map span
\[
f_!:=
\left(
X\xleftarrow{1_X}X\xrightarrow{f}Y
\right)
\]
and the reverse map span
\[
f^\ast_{\mathrm{sp}}:=
\left(
Y\xleftarrow{f}X\xrightarrow{1_X}X
\right).
\]

\begin{theorem}[Universal form of diamonds]
\label{thm:universal-diamond-relative}
In a regular category, the diamond assignment is the functorial extension from
map spans and reverse map spans determined by
\[
\langle f_!\rangle=f^\ast,
\qquad
\langle f^\ast_{\mathrm{sp}}\rangle=\exists_f.
\]
Equivalently, for a span
\[
R=
\left(
X\xleftarrow{\ell}R\xrightarrow{r}Y
\right)
=
\ell^\ast_{\mathrm{sp}};r_!,
\]
one has
\[
\langle R\rangle=\exists_\ell r^\ast.
\]
\end{theorem}

\begin{proof}
For a map span \(f_!\), the formula gives
\[
\exists_{1_X}f^\ast=f^\ast.
\]
For a reverse map span \(f^\ast_{\mathrm{sp}}\), the formula gives
\[
\exists_f 1_X^\ast=\exists_f.
\]
Every span factors as \(\ell^\ast_{\mathrm{sp}};r_!\), so functoriality forces
\[
\langle R\rangle
=
\langle \ell^\ast_{\mathrm{sp}}\rangle\langle r_!\rangle
=
\exists_\ell r^\ast.
\]
\end{proof}

\section{Relative primitive diamonds}
\label{sec:relative-primitive-diamonds}

Fix a Mealy morphism
\[
\Phi:\mathbb A\rightsquigarrow\mathbb B
\]
and a right \(\mathbb B\)-action \(Y\). The relative primitive execution span is
\[
\mathbf m_{\Phi,Y}
=
\left(
S_{\Phi,Y}
\xleftarrow{s_{\Phi,Y}}
M_{\Phi,Y}
\xrightarrow{t_{\Phi,Y}}
S_{\Phi,Y}
\right).
\]

For a predicate
\[
\varphi\hookrightarrow S_{\Phi,Y},
\]
the primitive diamond is
\[
\langle \mathbf m_{\Phi,Y}\rangle\varphi
=
\exists_{s_{\Phi,Y}}(t_{\Phi,Y}^\ast\varphi).
\]

\begin{proposition}[External reading of relative primitive diamonds]
\label{prop:external-reading-relative-primitive-diamonds}
In \(\mathbf{Set}\),
\[
(x,y)\models \langle \mathbf m_{\Phi,Y}\rangle\varphi
\]
if and only if there exists an \(A\)-arrow
\[
a:u\to p(x)
\]
such that
\[
\left(
\delta(a,x),
\sigma(\omega(a,x),y)
\right)
\models\varphi.
\]
\end{proposition}

\begin{proof}
This is the elementwise reading of
\[
\exists_{s_{\Phi,Y}}(t_{\Phi,Y}^\ast\varphi).
\]
A witness is precisely a primitive execution
\[
(a,x,y)\in M_{\Phi,Y}
\]
with source \((x,y)\), and the target of such an execution is
\[
\left(
\delta(a,x),
\sigma(\omega(a,x),y)
\right).
\]
\end{proof}

\begin{corollary}[Terminal primitive diamond]
\label{cor:terminal-primitive-diamond}
For \(Y=\mathbf 1_{\mathbb B}\), the primitive diamond reduces to the trace
diamond on \(P\):
\[
x\models\langle \mathbf m_\Phi\rangle\varphi
\]
if and only if there exists an admissible \(A\)-arrow \(a\) such that
\[
\delta(a,x)\models\varphi.
\]
The emitted output \(\omega(a,x)\) is available as a label through
\(\Omega_\Phi\), but it is not consumed by a nontrivial downstream action.
\end{corollary}

\section{Tests and relative state predicates}
\label{sec:tests-relative-state-predicates}

Let \(\mathcal E\) be regular. Fix \(\Phi\) and \(Y\), and write
\[
S:=S_{\Phi,Y}.
\]

\begin{definition}[Test]
\label{def:relative-test}
For a state predicate
\[
\theta:\Theta\hookrightarrow S,
\]
define the corresponding test span
\[
\theta?
=
\left(
S\xleftarrow{\theta}\Theta\xrightarrow{\theta}S
\right).
\]
\end{definition}

\begin{proposition}[Diamond law for tests]
\label{prop:relative-test-diamond-law}
For every predicate
\[
\varphi\hookrightarrow S,
\]
one has
\[
\langle\theta?\rangle\varphi
=
\theta\wedge\varphi.
\]
\end{proposition}

\begin{proof}
By definition,
\[
\langle\theta?\rangle\varphi
=
\exists_\theta(\theta^\ast\varphi).
\]
Since \(\theta\) is monic, existential image along \(\theta\) identifies
subobjects of \(\Theta\) with subobjects of \(S\) lying below \(\theta\). The
image of \(\theta^\ast\varphi\) is therefore
\[
\theta\wedge\varphi.
\]
\end{proof}

\begin{remark}
\label{rem:relative-tests-interpretation}
In the relative theory, tests are predicates on the coupled state object
\[
S_{\Phi,Y}=P\times_{B_0}Y.
\]
Thus a test may mention the Mealy state, the downstream \(Y\)-state, or any
relation between them. In the terminal trace case, tests reduce to predicates
on \(P\).
\end{remark}

\section{Input, output, and downstream transition restrictions}
\label{sec:input-output-downstream-restrictions}

The relative primitive execution object
\[
M_{\Phi,Y}
\]
supports several natural classes of subprograms.

Let
\[
\iota_A:M_{\Phi,Y}\to A_1
\]
be the input projection
\[
(a,x,y)\mapsto a.
\]
Let
\[
\iota_B:M_{\Phi,Y}\to B_1
\]
be the raw output map
\[
(a,x,y)\mapsto \omega(a,x).
\]
Let
\[
\lambda_{\Phi,Y}:M_{\Phi,Y}\to
\left(\int_{\mathbb B}Y\right)_1
\]
be the arrow map of the downstream comparison functor:
\[
(a,x,y)\mapsto(\omega(a,x),y).
\]

\begin{definition}[Restricted relative primitive subprogram]
\label{def:restricted-relative-subprogram}
Let
\[
R\hookrightarrow A_1
\]
be an input predicate, let
\[
O\hookrightarrow B_1
\]
be a raw output predicate, and let
\[
T\hookrightarrow \left(\int_{\mathbb B}Y\right)_1
\]
be a downstream transition predicate. Define
\[
R^\sharp:=\iota_A^\ast R,
\qquad
O^\sharp:=\iota_B^\ast O,
\qquad
T^\sharp:=\lambda_{\Phi,Y}^\ast T.
\]
Their simultaneous restriction is
\[
(R/O/T)^\sharp
:=
R^\sharp\wedge O^\sharp\wedge T^\sharp
\hookrightarrow M_{\Phi,Y}.
\]
The associated restricted primitive span is
\[
S_{\Phi,Y}
\xleftarrow{s i}
(R/O/T)^\sharp
\xrightarrow{t i}
S_{\Phi,Y},
\]
where
\[
i:(R/O/T)^\sharp\hookrightarrow M_{\Phi,Y}
\]
is the inclusion. Its diamond is denoted
\[
\langle R/O/T\rangle\varphi.
\]
\end{definition}

\begin{proposition}[External reading of restricted relative diamonds]
\label{prop:external-reading-relative-restricted}
In \(\mathbf{Set}\),
\[
(x,y)\models \langle R/O/T\rangle\varphi
\]
holds precisely when there exists an admissible \(A\)-arrow
\[
a:u\to p(x)
\]
such that
\[
R(a),
\qquad
O(\omega(a,x)),
\qquad
T(\omega(a,x),y),
\]
and
\[
\left(
\delta(a,x),
\sigma(\omega(a,x),y)
\right)
\models\varphi.
\]
\end{proposition}

\begin{proof}
The restricted primitive span is obtained by restricting the primitive
execution object \(M_{\Phi,Y}\) to those triples \((a,x,y)\) satisfying the
three pulled-back predicates. The diamond says that there exists such a
restricted primitive execution whose target state satisfies \(\varphi\). The
target state is
\[
\left(
\delta(a,x),
\sigma(\omega(a,x),y)
\right).
\]
\end{proof}

\begin{corollary}[Terminal restricted primitives]
\label{cor:terminal-restricted-primitives}
For \(Y=\mathbf 1_{\mathbb B}\), the downstream action category is
\[
\mathbb B^{\mathrm{op}}.
\]
Thus downstream transition predicates reduce to predicates on output arrows
of \(\mathbb B\), and the relative restricted modalities recover the trace
modalities that restrict primitive executions by input and output labels.
\end{corollary}

\section{Coherent finite choice of relative subprograms}
\label{sec:coherent-choice-relative}

\begin{definition}[Coherent category]
\label{def:coherent-category-relative}
A \textbf{coherent category} is a regular category \(\mathcal E\) such that:
\begin{enumerate}
\item for every object \(X\), the poset \(\mathrm{Sub}(X)\) has finite joins;
\item for every morphism \(f:X\to Y\), the pullback map
\[
f^\ast:\mathrm{Sub}(Y)\to\mathrm{Sub}(X)
\]
preserves finite joins.
\end{enumerate}
\end{definition}

Assume \(\mathcal E\) is coherent. Let
\[
U,V\hookrightarrow M_{\Phi,Y}
\]
be subprograms of the relative primitive execution span. Their finite choice
is the join
\[
U\vee V\hookrightarrow M_{\Phi,Y}.
\]

For a subobject \(i_U:U\hookrightarrow M_{\Phi,Y}\), write
\[
\langle U\rangle\varphi
:=
\exists_{s i_U}\bigl((t i_U)^\ast\varphi\bigr).
\]
Equivalently, regarding \(U\) as a predicate on \(M_{\Phi,Y}\),
\[
\langle U\rangle\varphi
=
\exists_s\bigl(U\wedge t^\ast\varphi\bigr).
\]

\begin{proposition}[Finite choice law]
\label{prop:relative-choice-law}
For subprograms
\[
U,V\hookrightarrow M_{\Phi,Y}
\]
in a coherent category,
\[
\langle U\vee V\rangle\varphi
=
\langle U\rangle\varphi
\vee
\langle V\rangle\varphi.
\]
\end{proposition}

\begin{proof}
Regard \(U\) and \(V\) as predicates on \(M_{\Phi,Y}\). Then
\[
\langle U\vee V\rangle\varphi
=
\exists_s\bigl((U\vee V)\wedge t^\ast\varphi\bigr).
\]
In a coherent category, finite joins distribute over finite meets. Hence
\[
(U\vee V)\wedge t^\ast\varphi
=
(U\wedge t^\ast\varphi)
\vee
(V\wedge t^\ast\varphi).
\]
Since \(\exists_s\) is a left adjoint, it preserves finite joins. Therefore
\[
\langle U\vee V\rangle\varphi
=
\langle U\rangle\varphi
\vee
\langle V\rangle\varphi.
\]
\end{proof}

\begin{remark}
\label{rem:relative-choice-not-arbitrary-coproduct}
This is finite choice of subprograms inside a fixed relative primitive
execution span. It is not arbitrary coproduct choice of spans. The operation is
the join of subobjects of \(M_{\Phi,Y}\).
\end{remark}

\section{Boxes and intuitionistic tests}
\label{sec:relative-boxes}

\begin{definition}[First-order Heyting category]
\label{def:first-order-heyting-relative}
A \textbf{first-order Heyting category} is a coherent category \(\mathcal E\)
such that for every morphism
\[
f:X\to Y,
\]
the pullback functor
\[
f^\ast:\mathrm{Sub}(Y)\to\mathrm{Sub}(X)
\]
has a right adjoint
\[
\forall_f:\mathrm{Sub}(X)\to\mathrm{Sub}(Y),
\]
and these right adjoints satisfy Beck--Chevalley for pullback squares.
Thus, for every \(f\), there is an adjoint triple
\[
\exists_f\dashv f^\ast\dashv \forall_f.
\]
\end{definition}

Assume \(\mathcal E\) is a first-order Heyting category.

\begin{definition}[Box modality of a span]
\label{def:span-box-relative}
For a span
\[
R=
\left(
X\xleftarrow{\ell}R\xrightarrow{r}Y
\right),
\]
define
\[
[R]:\mathrm{Sub}(Y)\to\mathrm{Sub}(X)
\]
by
\[
[R]\varphi:=\forall_\ell(r^\ast\varphi).
\]
\end{definition}

\begin{proposition}[Functoriality of boxes]
\label{prop:box-functoriality-relative}
For spans
\[
R:X\nrightarrow Y
\qquad\text{and}\qquad
S:Y\nrightarrow Z
\]
in a first-order Heyting category,
\[
[1_X]=\mathrm{id}_{\mathrm{Sub}(X)}
\]
and
\[
[R;S]=[R]\circ[S].
\]
\end{proposition}

\begin{proof}
The identity law is immediate. For composition, write
\[
R=(X\xleftarrow{\ell}R\xrightarrow{r}Y),
\qquad
S=(Y\xleftarrow{m}S\xrightarrow{n}Z).
\]
The composite has apex \(R\times_Y S\), so
\[
[R;S]\varphi
=
\forall_{\ell\pi_R}\bigl((n\pi_S)^\ast\varphi\bigr).
\]
Using functoriality of universal image and Beck--Chevalley for the pullback
square defining \(R\times_Y S\), this becomes
\[
\forall_\ell r^\ast\forall_m n^\ast\varphi
=
[R]([S]\varphi).
\]
\end{proof}

\begin{proposition}[Box law for tests]
\label{prop:relative-box-test-law}
For a test
\[
\theta?
=
(S\xleftarrow{\theta}\Theta\xrightarrow{\theta}S)
\]
in a first-order Heyting category,
\[
[\theta?]\varphi
=
\theta\Rightarrow\varphi.
\]
\end{proposition}

\begin{proof}
By definition,
\[
[\theta?]\varphi
=
\forall_\theta(\theta^\ast\varphi).
\]
For any predicate \(\psi\hookrightarrow S\),
\[
\psi\le \forall_\theta(\theta^\ast\varphi)
\]
if and only if
\[
\theta^\ast\psi\le \theta^\ast\varphi
\]
by the adjunction \(\theta^\ast\dashv\forall_\theta\). Since \(\theta\) is
monic, this is equivalent to
\[
\psi\wedge\theta\le\varphi.
\]
By the defining property of Heyting implication, this is equivalent to
\[
\psi\le \theta\Rightarrow\varphi.
\]
\end{proof}

\begin{theorem}[Universal form of boxes]
\label{thm:universal-box-relative}
In a first-order Heyting category, the box assignment is the functorial
extension from map spans and reverse map spans determined by
\[
[f_!]=f^\ast,
\qquad
[f^\ast_{\mathrm{sp}}]=\forall_f.
\]
Equivalently, for a span
\[
R=
\left(
X\xleftarrow{\ell}R\xrightarrow{r}Y
\right)
=
\ell^\ast_{\mathrm{sp}};r_!,
\]
one has
\[
[R]=\forall_\ell r^\ast.
\]
\end{theorem}

\begin{proof}
For a map span \(f_!\), the formula gives
\[
\forall_{1_X}f^\ast=f^\ast.
\]
For a reverse map span \(f^\ast_{\mathrm{sp}}\), the formula gives
\[
\forall_f 1_X^\ast=\forall_f.
\]
Every span factors as \(\ell^\ast_{\mathrm{sp}};r_!\), so functoriality forces
\[
[R]
=
[\ell^\ast_{\mathrm{sp}}][r_!]
=
\forall_\ell r^\ast.
\]
\end{proof}

\begin{remark}
\label{rem:box-not-boolean-dual-relative}
Boxes are not defined as Boolean duals of diamonds. Structurally,
\[
\langle R\rangle\varphi=\exists_\ell(r^\ast\varphi)
\]
uses the left adjoint to pullback along the source leg, while
\[
[R]\varphi=\forall_\ell(r^\ast\varphi)
\]
uses the right adjoint. In Boolean settings one may recover boxes from
diamonds by duality, but the definition above is meaningful intuitionistically.
\end{remark}

\section{Change of downstream model and modal stability}
\label{sec:downstream-change-modal-stability}

Let
\[
f:Y\to Z
\]
be a morphism of right \(\mathbb B\)-actions. It induces a map of relative
state objects
\[
h_f:S_{\Phi,Y}\to S_{\Phi,Z},
\qquad
h_f(x,y)=(x,f(y)).
\]
It also induces a map of primitive execution objects
\[
k_f:M_{\Phi,Y}\to M_{\Phi,Z},
\qquad
k_f(a,x,y)=(a,x,f(y)).
\]

Let
\[
V\hookrightarrow M_{\Phi,Z}
\]
be a subprogram. Its pullback along \(k_f\) is a subprogram
\[
k_f^\ast V\hookrightarrow M_{\Phi,Y}.
\]

\begin{proposition}[Diamond stability under change of downstream model]
\label{prop:diamond-stability-downstream-change}
Assume \(\mathcal E\) is regular. For every predicate
\[
\psi\hookrightarrow S_{\Phi,Z},
\]
one has
\[
h_f^\ast\langle V\rangle\psi
=
\langle k_f^\ast V\rangle h_f^\ast\psi.
\]
\end{proposition}

\begin{proof}
The square
\[
\begin{tikzcd}
M_{\Phi,Y} \arrow[r,"k_f"] \arrow[d,"s_Y"']
&
M_{\Phi,Z} \arrow[d,"s_Z"]
\\
S_{\Phi,Y} \arrow[r,"h_f"']
&
S_{\Phi,Z}
\end{tikzcd}
\]
is a pullback, since both primitive execution objects are obtained by pulling
back \(A_1\) along the corresponding state object over \(A_0\). Pulling back
the subobject \(V\hookrightarrow M_{\Phi,Z}\) gives a pullback square
\[
\begin{tikzcd}
k_f^\ast V \arrow[r] \arrow[d,"s_Y"']
&
V \arrow[d,"s_Z"]
\\
S_{\Phi,Y} \arrow[r,"h_f"']
&
S_{\Phi,Z}.
\end{tikzcd}
\]
Beck--Chevalley for regular images gives the corresponding equality for
existential images along the source legs. The target maps commute with
\(h_f\) because \(f\) is equivariant, so pulling back \(t_Z^\ast\psi\) along
the top map is the same as pulling back \(h_f^\ast\psi\) along \(t_Y\). This
gives the displayed identity.
\end{proof}

\begin{proposition}[Box stability under change of downstream model]
\label{prop:box-stability-downstream-change}
Assume \(\mathcal E\) is a first-order Heyting category. For every predicate
\[
\psi\hookrightarrow S_{\Phi,Z},
\]
one has
\[
h_f^\ast[V]\psi
=
[k_f^\ast V]h_f^\ast\psi.
\]
\end{proposition}

\begin{proof}
The proof is the same as that of
Proposition~\ref{prop:diamond-stability-downstream-change}, using
Beck--Chevalley for the right adjoints \(\forall\) instead of for the left
adjoints \(\exists\).
\end{proof}

\begin{remark}
\label{rem:modal-stability-interpretation}
Changing the downstream model commutes with modal interpretation, provided
subprograms are pulled back along the induced program functor. This is the
logical counterpart of the functoriality
\[
Y\mapsto \mathsf{Prog}_{\Phi,Y}.
\]
\end{remark}

\section{Finite sequential iteration}
\label{sec:relative-finite-sequential-iteration}

Fix \(\Phi\) and \(Y\), and write
\[
S:=S_{\Phi,Y}.
\]
Let
\[
U:S\nrightarrow S
\]
be an endoprogram, usually a subspan of the relative primitive execution span
\[
\mathbf m_{\Phi,Y}.
\]

Define
\[
U^0:=1_S,
\qquad
U^{n+1}:=U;U^n.
\]
When necessary, write
\[
U^n=
\left(
S\xleftarrow{\ell_n}U_n\xrightarrow{r_n}S
\right),
\]
where \(U_n\) denotes the apex of the \(n\)-fold composite.

\begin{proposition}[Finite sequential laws]
\label{prop:relative-finite-sequential-laws}
After the canonical associativity and unit identifications of spans,
\[
U^{m+n}\cong U^m;U^n.
\]
Consequently, in a regular category,
\[
\langle U^{m+n}\rangle
=
\langle U^m\rangle\circ\langle U^n\rangle,
\]
and in a first-order Heyting category,
\[
[U^{m+n}]
=
[U^m]\circ[U^n].
\]
\end{proposition}

\begin{proof}
The span isomorphism follows by induction on \(m\), using the associativity and
unit constraints of \(\mathbf{Span}(\mathcal E)\). The modal identities follow
from Proposition~\ref{prop:diamond-functoriality-relative} and
Proposition~\ref{prop:box-functoriality-relative}.
\end{proof}

\section{Relative path star and its recursive equation}
\label{sec:relative-path-star-recursive-equation}

To form the span of all finite paths of an endoprogram, we require stable
countable coproducts.

\begin{definition}[Path-star span]
\label{def:relative-path-star}
Suppose \(\mathcal E\) has countable coproducts stable under pullback. Let
\[
U:S\nrightarrow S
\]
be an endospan, and write
\[
U^n=
(S\xleftarrow{\ell_n}U_n\xrightarrow{r_n}S).
\]
The \textbf{finite-path star} of \(U\) is the endospan
\[
U^\star_{\mathrm{path}}:S\nrightarrow S
\]
whose apex is
\[
\coprod_{n\ge0}U_n
\]
and whose legs are induced by the legs of the powers:
\[
S
\xleftarrow{[\ell_n]_{n\ge0}}
\coprod_{n\ge0}U_n
\xrightarrow{[r_n]_{n\ge0}}
S.
\]
\end{definition}

\begin{proposition}[Path-star equation]
\label{prop:relative-path-star-equation}
There is a canonical isomorphism of endospans
\[
U^\star_{\mathrm{path}}
\cong
1_S\oplus U;U^\star_{\mathrm{path}},
\]
where \(\oplus\) denotes coproduct of span apexes.
\end{proposition}

\begin{proof}
By definition, the apex of \(U^\star_{\mathrm{path}}\) is
\[
\coprod_{n\ge0}U_n.
\]
Separating the \(0\)-summand gives
\[
\coprod_{n\ge0}U_n
\cong
U_0+\coprod_{n\ge0}U_{n+1}.
\]
Since
\[
U^0=1_S
\]
and
\[
U^{n+1}=U;U^n,
\]
the second summand is
\[
\coprod_{n\ge0}\operatorname{apex}(U;U^n).
\]
Because countable coproducts are stable under pullback, composing \(U\) with
the coproduct span \(\coprod_{n\ge0}U^n\) gives the coproduct of the composites:
\[
U;\left(\coprod_{n\ge0}U^n\right)
\cong
\coprod_{n\ge0}(U;U^n).
\]
Thus
\[
U^\star_{\mathrm{path}}
\cong
1_S\oplus U;U^\star_{\mathrm{path}}.
\]
\end{proof}

\section{Relative path star as a free endospan monoid}
\label{sec:relative-path-star-free-monoid}

\begin{theorem}[Path star as free endospan monoid]
\label{thm:relative-path-star-free-monoid}
Let \(\mathcal E\) have pullbacks and countable coproducts stable under
pullback. For every endospan
\[
U:S\nrightarrow S,
\]
the span
\[
U^\star_{\mathrm{path}}
=
\coprod_{n\ge0}U^n
\]
is the free monoid on \(U\) in the monoidal category
\[
\operatorname{EndSpan}_{\mathcal E}(S).
\]
\end{theorem}

\begin{proof}
Because countable coproducts are stable under pullback, composition with a
fixed span preserves countable coproducts of apexes on either side. Thus, for
every fixed span \(R\), both functors
\[
R;(-)
\qquad\text{and}\qquad
(-);R
\]
preserve countable coproducts of spans with fixed endpoints. Hence
\(\operatorname{EndSpan}_{\mathcal E}(S)\) has countable coproducts, computed
by coproducts of apexes, and the monoidal product preserves these coproducts
separately in each variable.

Define
\[
U^\star_{\mathrm{path}}:=\coprod_{n\ge0}U^n.
\]
The unit is the inclusion of the \(0\)-summand
\[
U^0=1_S\to U^\star_{\mathrm{path}}.
\]
For multiplication, use distributivity of composition over coproducts:
\[
U^\star_{\mathrm{path}};U^\star_{\mathrm{path}}
\cong
\coprod_{m,n\ge0}(U^m;U^n).
\]
By Proposition~\ref{prop:relative-finite-sequential-laws}, each summand is
canonically isomorphic to \(U^{m+n}\), followed by the coproduct inclusion
\[
U^{m+n}\to\coprod_{k\ge0}U^k.
\]
This defines the multiplication. Associativity and unit laws follow from
associativity and unit laws of span composition and addition of natural
numbers.

Given a monoid object \(M\) in \(\operatorname{EndSpan}_{\mathcal E}(S)\) and
a map of spans \(i:U\to M\), define maps \(i_n:U^n\to M\) by
\[
i_0:=\eta,
\qquad
i_1:=i,
\]
and
\[
i_{n+1}:=\mu\circ(i;i_n),
\]
where \(\eta\) and \(\mu\) are the unit and multiplication of \(M\). The
coproduct universal property gives a unique map
\[
\overline i:U^\star_{\mathrm{path}}\to M
\]
restricting to \(i_n\) on \(U^n\). The monoid morphism laws follow by induction
on summands, and uniqueness follows because every monoid morphism extending
\(i\) must have these restrictions.
\end{proof}

\begin{remark}
\label{rem:relative-two-monoids}
For a relative program category \(\mathsf{Prog}_{\Phi,Y}\), there are two
different monoids in view. The primitive execution span
\[
\mathbf m_{\Phi,Y}
\]
is the ambient execution monoid forced by the Eilenberg--Moore action
\(\Phi^\ast Y\). If
\[
U\hookrightarrow \mathbf m_{\Phi,Y}
\]
is a chosen subprogram, then
\[
U^\star_{\mathrm{path}}
\]
is the freely generated path monoid of that subprogram.
\end{remark}

\section{Countable predicate structure and fixed points}
\label{sec:relative-fixed-points}

\begin{definition}[\(\omega\)-iterative first-order Heyting category]
\label{def:omega-iterative-relative}
An \(\omega\)-iterative first-order Heyting category is a first-order Heyting
category \(\mathcal E\) such that:
\begin{enumerate}
\item \(\mathcal E\) has countable coproducts;
\item countable coproducts are stable under pullback;
\item countable coproducts are subobject-extensive, in the sense that for every
countable family \((X_n)_{n\ge0}\), pullback along the coproduct injections
induces an isomorphism
\[
\mathrm{Sub}\left(\coprod_{n\ge0}X_n\right)
\cong
\prod_{n\ge0}\mathrm{Sub}(X_n);
\]
\item for every object \(X\), the fibre \(\mathrm{Sub}(X)\) is a complete
lattice.
\end{enumerate}
\end{definition}

\begin{remark}[Strength of the \(\omega\)-iterative hypothesis]
\label{rem:omega-iterative-hypothesis-strength-relative}
For the specific Kleene-chain comparisons below, it is enough to have the
countable joins, countable meets, and continuity properties that occur in the
displayed formulas. The complete-fibre formulation makes the fixed-point
terminology unambiguous and keeps the statement stable under later extensions
of the logic.
\end{remark}

\begin{example}[Grothendieck topoi]
\label{ex:grothendieck-topos-relative}
Every Grothendieck topos is an \(\omega\)-iterative first-order Heyting
category. Consequently, the relative path-star and fixed-point constructions
below apply in every Grothendieck topos, for Mealy data and downstream actions
internal to that topos.
\end{example}

Assume for the rest of this section that \(\mathcal E\) is
\(\omega\)-iterative first-order Heyting.

Let
\[
U:S\nrightarrow S
\]
be an endoprogram, and let
\[
\varphi\in\mathrm{Sub}(S).
\]

\subsection{Diamond fixed points}

Define
\[
F^\Diamond_{U,\varphi}:\mathrm{Sub}(S)\to\mathrm{Sub}(S)
\]
by
\[
F^\Diamond_{U,\varphi}(X)
:=
\varphi\vee\langle U\rangle X.
\]
Since
\[
r^\ast\dashv \forall_r,
\]
the pullback \(r^\ast\) preserves all existing joins, and since
\[
\exists_\ell\dashv \ell^\ast,
\]
the existential image \(\exists_\ell\) preserves all existing joins. Hence
\[
\langle U\rangle=\exists_\ell r^\ast
\]
preserves countable joins. Therefore \(F^\Diamond_{U,\varphi}\) is
\(\omega\)-continuous for the ascending Kleene chain from \(\bot\). Define
\[
\mu F^\Diamond_{U,\varphi}
:=
\bigvee_{n\ge0}(F^\Diamond_{U,\varphi})^n(\bot).
\]

\begin{lemma}[Diamond approximants]
\label{lem:relative-diamond-approximants}
For every \(n\ge0\),
\[
(F^\Diamond_{U,\varphi})^n(\bot)
=
\bigvee_{0\le k<n}\langle U^k\rangle\varphi.
\]
\end{lemma}

\begin{proof}
The case \(n=0\) is the empty join. The induction step uses
\[
F^\Diamond_{U,\varphi}(X)=\varphi\vee\langle U\rangle X,
\]
preservation of joins by \(\langle U\rangle\), and functoriality
\[
\langle U\rangle\langle U^k\rangle
=
\langle U;U^k\rangle
=
\langle U^{k+1}\rangle.
\]
\end{proof}

\begin{theorem}[Diamond least fixed point]
\label{thm:relative-diamond-lfp}
In \(\mathrm{Sub}(S)\),
\[
\mu X.\bigl(\varphi\vee\langle U\rangle X\bigr)
=
\bigvee_{n\ge0}\langle U^n\rangle\varphi.
\]
This predicate is the least pre-fixed point of
\[
X\mapsto \varphi\vee\langle U\rangle X.
\]
\end{theorem}

\begin{proof}
The displayed formula follows from
Lemma~\ref{lem:relative-diamond-approximants}. Let
\[
\chi:=\bigvee_{n\ge0}\langle U^n\rangle\varphi.
\]
Then
\[
\begin{aligned}
\varphi\vee\langle U\rangle\chi
&=
\langle U^0\rangle\varphi
\vee
\bigvee_{n\ge0}\langle U\rangle\langle U^n\rangle\varphi
\\
&=
\langle U^0\rangle\varphi
\vee
\bigvee_{n\ge0}\langle U^{n+1}\rangle\varphi
\\
&=
\chi.
\end{aligned}
\]
Thus \(\chi\) is a fixed point. If
\[
\varphi\vee\langle U\rangle\psi\le \psi,
\]
then \(\varphi\le\psi\) and \(\langle U\rangle\psi\le\psi\). By induction on
\(n\),
\[
\langle U^n\rangle\varphi\le\psi
\]
for all \(n\). Therefore \(\chi\le\psi\), so \(\chi\) is the least pre-fixed
point.
\end{proof}

\begin{corollary}[Diamond induction principle]
\label{cor:relative-diamond-induction}
If
\[
\varphi\vee\langle U\rangle\psi\le\psi,
\]
then
\[
\mu X.\bigl(\varphi\vee\langle U\rangle X\bigr)\le\psi.
\]
\end{corollary}

\begin{proof}
This is exactly the least pre-fixed point property from
Theorem~\ref{thm:relative-diamond-lfp}.
\end{proof}

\subsection{Box fixed points}

Define
\[
F^\Box_{U,\varphi}:\mathrm{Sub}(S)\to\mathrm{Sub}(S)
\]
by
\[
F^\Box_{U,\varphi}(X)
:=
\varphi\wedge[U]X.
\]
Since
\[
\exists_r\dashv r^\ast,
\]
the pullback \(r^\ast\) preserves all existing meets, and since
\[
\ell^\ast\dashv \forall_\ell,
\]
the universal image \(\forall_\ell\) preserves all existing meets. Hence
\[
[U]=\forall_\ell r^\ast
\]
preserves countable meets. Therefore \(F^\Box_{U,\varphi}\) preserves the
relevant decreasing \(\omega\)-chains. Define
\[
\nu F^\Box_{U,\varphi}
:=
\bigwedge_{n\ge0}(F^\Box_{U,\varphi})^n(\top).
\]

\begin{lemma}[Box approximants]
\label{lem:relative-box-approximants}
For every \(n\ge0\),
\[
(F^\Box_{U,\varphi})^n(\top)
=
\bigwedge_{0\le k<n}[U^k]\varphi.
\]
\end{lemma}

\begin{proof}
The case \(n=0\) is the empty meet. The induction step uses preservation of
meets by \([U]\) and functoriality
\[
[U][U^k]=[U;U^k]=[U^{k+1}].
\]
\end{proof}

\begin{theorem}[Box greatest fixed point]
\label{thm:relative-box-gfp}
In \(\mathrm{Sub}(S)\),
\[
\nu X.\bigl(\varphi\wedge[U]X\bigr)
=
\bigwedge_{n\ge0}[U^n]\varphi.
\]
This predicate is the greatest post-fixed point of
\[
X\mapsto \varphi\wedge[U]X.
\]
\end{theorem}

\begin{proof}
The displayed formula follows from
Lemma~\ref{lem:relative-box-approximants}. Let
\[
\chi:=\bigwedge_{n\ge0}[U^n]\varphi.
\]
Then
\[
\begin{aligned}
\varphi\wedge[U]\chi
&=
[U^0]\varphi
\wedge
\bigwedge_{n\ge0}[U][U^n]\varphi
\\
&=
[U^0]\varphi
\wedge
\bigwedge_{n\ge0}[U^{n+1}]\varphi
\\
&=
\chi.
\end{aligned}
\]
Thus \(\chi\) is a fixed point. If
\[
\psi\le \varphi\wedge[U]\psi,
\]
then \(\psi\le\varphi\) and \(\psi\le[U]\psi\). By induction on \(n\),
\[
\psi\le[U^n]\varphi
\]
for all \(n\). Hence
\[
\psi\le\chi.
\]
Thus \(\chi\) is the greatest post-fixed point.
\end{proof}

\begin{corollary}[Box coinduction principle]
\label{cor:relative-box-coinduction}
If
\[
\psi\le \varphi\wedge[U]\psi,
\]
then
\[
\psi\le
\nu X.\bigl(\varphi\wedge[U]X\bigr).
\]
\end{corollary}

\begin{proof}
This is exactly the greatest post-fixed point property from
Theorem~\ref{thm:relative-box-gfp}.
\end{proof}

\section{Path semantics and fixed-point semantics}
\label{sec:relative-path-fixed-point-comparison}

Let \(\mathcal E\) be an \(\omega\)-iterative first-order Heyting category.
Let
\[
U:S\nrightarrow S
\]
be an endoprogram, and let
\[
U^\star_{\mathrm{path}}
=
\coprod_{n\ge0}U^n
\]
be its finite-path star.

\begin{proposition}[Modalities of countable coproduct spans]
\label{prop:relative-coproduct-modalities}
Let \((R_n)_{n\ge0}\) be a countable family of spans
\[
R_n:X\nrightarrow Y.
\]
Then
\[
\left\langle\coprod_{n\ge0}R_n\right\rangle\varphi
=
\bigvee_{n\ge0}\langle R_n\rangle\varphi
\]
and
\[
\left[\coprod_{n\ge0}R_n\right]\varphi
=
\bigwedge_{n\ge0}[R_n]\varphi.
\]
\end{proposition}

\begin{proof}
The proof uses subobject-extensivity of countable coproducts. A subobject of
\[
\coprod_n R_n
\]
is determined by its restrictions to the summands. Existential image along the
coproduct leg is therefore the join of the existential images along the
summand legs, giving the diamond formula. Dually, the defining adjunction for
\(\forall\) shows that universal image along the coproduct leg is the meet of
the universal images along the summand legs, giving the box formula.
\end{proof}

\begin{theorem}[Path/fixed-point comparison]
\label{thm:relative-path-fixed-point-comparison}
For every predicate
\[
\varphi\in\mathrm{Sub}(S),
\]
one has
\[
\langle U^\star_{\mathrm{path}}\rangle\varphi
=
\mu X.\bigl(\varphi\vee\langle U\rangle X\bigr),
\]
and
\[
[U^\star_{\mathrm{path}}]\varphi
=
\nu X.\bigl(\varphi\wedge[U]X\bigr).
\]
\end{theorem}

\begin{proof}
By Proposition~\ref{prop:relative-coproduct-modalities},
\[
\langle U^\star_{\mathrm{path}}\rangle\varphi
=
\bigvee_{n\ge0}\langle U^n\rangle\varphi.
\]
By Theorem~\ref{thm:relative-diamond-lfp}, this is
\[
\mu X.\bigl(\varphi\vee\langle U\rangle X\bigr).
\]
The box statement is identical, using
Proposition~\ref{prop:relative-coproduct-modalities} and
Theorem~\ref{thm:relative-box-gfp}.
\end{proof}

\begin{corollary}[Grothendieck topos case]
\label{cor:grothendieck-topos-relative-path-fixed-point}
If \(\mathcal E\) is a Grothendieck topos, then for every Mealy morphism
\[
\Phi:\mathbb A\rightsquigarrow\mathbb B,
\]
every right \(\mathbb B\)-action \(Y\), every relative endoprogram
\[
U:S_{\Phi,Y}\nrightarrow S_{\Phi,Y},
\]
and every predicate
\[
\varphi\in\mathrm{Sub}(S_{\Phi,Y}),
\]
the path/fixed-point comparison of
Theorem~\ref{thm:relative-path-fixed-point-comparison} applies.
\end{corollary}

\begin{proof}
By Example~\ref{ex:grothendieck-topos-relative}, every Grothendieck topos is
an \(\omega\)-iterative first-order Heyting category.
\end{proof}

\begin{remark}
\label{rem:relative-path-fixed-point-meaning}
The path-star span is witnessful: it remembers a finite execution path as a
path witness. The fixed-point modalities live in the predicate fibre
\[
\mathrm{Sub}(S).
\]
The theorem says that, under the stated hypotheses, witnessful finite-path
semantics and fibrewise fixed-point semantics agree after applying the span
predicate semantics.
\end{remark}

\section{Relative folded trace semantics}
\label{sec:relative-folded-trace-semantics}

Let
\[
i:U\hookrightarrow M_{\Phi,Y}
\]
be a subprogram of the relative primitive execution span
\[
\mathbf m_{\Phi,Y}.
\]
Equivalently, \(i\) is a map of spans
\[
U\to\mathbf m_{\Phi,Y}.
\]

Since
\[
\mathbf m_{\Phi,Y}
\]
is a monoid object in
\[
\operatorname{EndSpan}_{\mathcal E}(S_{\Phi,Y}),
\]
the universal property of the free path monoid gives a canonical folded trace
map.

\begin{definition}[Relative folded trace map]
\label{def:relative-folded-trace-map}
The \textbf{relative folded trace map} of a subprogram
\[
i:U\hookrightarrow \mathbf m_{\Phi,Y}
\]
is the unique monoid morphism
\[
\operatorname{fold}_{U,Y}:
U^\star_{\mathrm{path}}\to\mathbf m_{\Phi,Y}
\]
whose restriction to the \(1\)-step summand \(U\) is \(i\).
\end{definition}

Let
\[
\mu_{\Phi,Y}:
M_{\Phi,Y}\times_{t_{\Phi,Y},S_{\Phi,Y},s_{\Phi,Y}}M_{\Phi,Y}
\to
M_{\Phi,Y}
\]
denote the multiplication map of the primitive execution monoid
\(\mathbf m_{\Phi,Y}\).

Equivalently, \(\operatorname{fold}_{U,Y}\) is assembled from maps
\[
\operatorname{fold}_n:U_n\to M_{\Phi,Y}
\]
defined recursively. For \(n=0\), set
\[
\operatorname{fold}_0:=e_{\Phi,Y}:S_{\Phi,Y}\to M_{\Phi,Y}.
\]
For \(n=1\), set
\[
\operatorname{fold}_1:=i:U\to M_{\Phi,Y}.
\]
For \(n\ge1\), the apex of
\[
U^{n+1}=U;U^n
\]
is
\[
U_{n+1}
=
U\times_{r_U,S_{\Phi,Y},\ell_n}U_n,
\]
where
\[
U=
(S_{\Phi,Y}\xleftarrow{\ell_U}U\xrightarrow{r_U}S_{\Phi,Y})
\]
and
\[
U^n=
(S_{\Phi,Y}\xleftarrow{\ell_n}U_n\xrightarrow{r_n}S_{\Phi,Y}).
\]
Define
\[
\operatorname{fold}_{n+1}
:=
\mu_{\Phi,Y}\circ(i\times\operatorname{fold}_n),
\]
where
\[
i\times\operatorname{fold}_n:
U\times_{r_U,S_{\Phi,Y},\ell_n}U_n
\to
M_{\Phi,Y}\times_{t_{\Phi,Y},S_{\Phi,Y},s_{\Phi,Y}}M_{\Phi,Y}
\]
is the induced map on pullbacks.

\begin{proposition}[Fold maps are maps of spans]
\label{prop:relative-fold-maps-span-maps}
Each
\[
\operatorname{fold}_n:U_n\to M_{\Phi,Y}
\]
is a map of spans
\[
U^n\Rightarrow \mathbf m_{\Phi,Y}.
\]
Consequently, by the coproduct universal property, the maps
\[
\operatorname{fold}_n
\]
assemble into the unique monoid morphism
\[
\operatorname{fold}_{U,Y}:
U^\star_{\mathrm{path}}\to\mathbf m_{\Phi,Y}
\]
extending
\[
i:U\to\mathbf m_{\Phi,Y}.
\]
\end{proposition}

\begin{proof}
For \(n=0\), this is the unit map of the primitive execution monoid. For
\(n=1\), it is the inclusion of the subprogram. Assume the claim holds for
\(n\ge1\). Then
\[
i\times\operatorname{fold}_n
\]
is a map from the composite span \(U;U^n\) to the composite span
\(\mathbf m_{\Phi,Y};\mathbf m_{\Phi,Y}\). Composing with the multiplication
map
\[
\mu_{\Phi,Y}:
\mathbf m_{\Phi,Y};\mathbf m_{\Phi,Y}\to\mathbf m_{\Phi,Y}
\]
gives a map of spans
\[
U^{n+1}\Rightarrow \mathbf m_{\Phi,Y}.
\]
The final statement follows from the coproduct universal property and the free
monoid theorem.
\end{proof}

\begin{proposition}[Folded downstream transitions]
\label{prop:relative-folded-downstream-transitions}
The downstream comparison functor
\[
\Lambda_{\Phi,Y}:\mathsf{Prog}_{\Phi,Y}\to\int_{\mathbb B}Y
\]
sends the folded composite of a finite relative execution path to the
corresponding composite transition in the downstream action category
\[
\int_{\mathbb B}Y.
\]
\end{proposition}

\begin{proof}
The folded trace is a composite arrow in
\[
\mathsf{Prog}_{\Phi,Y}.
\]
Since \(\Lambda_{\Phi,Y}\) is an internal functor, it preserves identities and
composition. Hence it sends the folded composite to the composite of the
downstream transitions of the primitive steps.
\end{proof}

\begin{corollary}[Terminal folded outputs]
\label{cor:terminal-folded-outputs}
For \(Y=\mathbf 1_{\mathbb B}\), the functor
\[
\Lambda_{\Phi,Y}
\]
is the output labelling functor
\[
\Omega_\Phi:\mathsf{Prog}_\Phi\to\mathbb B^{\mathrm{op}}.
\]
Thus terminal folded traces recover folded output arrows in
\(\mathbb B^{\mathrm{op}}\).
\end{corollary}

\begin{remark}
\label{rem:path-versus-fold-relative}
Path semantics remembers the decomposition of an execution into primitive
steps. Folded trace semantics remembers only the composite arrow in the
relative program category. In output-sensitive or downstream-sensitive
logics, this distinction matters: requiring every primitive output to satisfy
a predicate is different from requiring only the folded downstream transition
to satisfy a predicate.
\end{remark}

\section{Classical and relational special cases}
\label{sec:classical-relational-special-cases}

\begin{example}[Classical Mealy cascades]
\label{ex:classical-mealy-cascades}
Let \(\mathcal E=\mathbf{Set}\), and suppose that \(\mathbb A\) and
\(\mathbb B\) each have one object. Then they are monoids, say
\[
\mathbb A=M,
\qquad
\mathbb B=N.
\]
A Mealy morphism consists of a set \(P\) and functions
\[
\delta:M\times P\to P,
\qquad
\omega:M\times P\to N
\]
satisfying the usual unit and cocycle laws, with the multiplication order
determined by the categorical convention of this chapter.

A right \(N\)-action is a set \(Y\) with an action
\[
N\times Y\to Y.
\]
The restricted \(M\)-action is on
\[
P\times Y
\]
and is given by
\[
m\cdot(x,y)
=
\left(
\delta(m,x),
\omega(m,x)\cdot y
\right).
\]
Thus the relative construction is exactly the automata-theoretic cascade of a
Mealy transducer followed by an \(N\)-automaton.

The terminal case has \(Y=1\). Then the state space is just \(P\), and the
output is recorded rather than consumed by a nontrivial downstream automaton.
\end{example}

\begin{example}[Words versus letter steps]
\label{ex:words-versus-letter-steps-relative}
If
\[
\mathbb A=I^\ast
\]
is the one-object category associated to the free monoid of input words, then
a primitive execution of the relative program category consumes a whole word,
not necessarily a single input letter. The usual letter-step presentation of a
classical Mealy machine is recovered either by restricting the primitive
execution object to the subpredicate of length-one words, or by starting with a
generator-level transition-output map and extending it to \(I^\ast\).
\end{example}

\begin{example}[Relational propositional dynamic logic]
\label{ex:relative-relational-pdl}
Let \(\mathcal E=\mathbf{Set}\) and pass from spans to their extensional
binary relations. A span
\[
S\xleftarrow{\ell}R\xrightarrow{r}S
\]
determines the relation
\[
x\,\overline R\,y
\]
if and only if there exists \(\rho\in R\) such that
\[
\ell(\rho)=x,
\qquad
r(\rho)=y.
\]
The diamond becomes the usual relational dynamic-logic modality:
\[
x\models\langle R\rangle\varphi
\]
if and only if there exists \(y\) such that
\[
x\,\overline R\,y
\]
and
\[
y\models\varphi.
\]

In the relative Mealy setting, this relational semantics is applied to the
coupled state space
\[
S_{\Phi,Y}=P\times_{B_0}Y.
\]
Thus ordinary relational PDL is recovered as the extensional quotient of the
relative span semantics.
\end{example}

\begin{example}[Relational Kleene algebra with tests]
\label{ex:relative-relational-kat}
In \(\mathbf{Set}\), after quotienting spans by extensional relations,
endoprograms on a set \(S\) form the relation algebra
\[
\mathcal P(S\times S).
\]
Sequential composition is relational composition, finite choice is union,
tests are subidentity relations, and path star is reflexive-transitive closure.

For the relative Mealy state space
\[
S_{\Phi,Y}=P\times_{B_0}Y,
\]
this gives a relational Kleene algebra with tests of the coupled system. The
terminal trace case is obtained by setting
\[
Y=\mathbf 1_{\mathbb B}.
\]
\end{example}

\section{Progressive consolidation}
\label{sec:relative-progressive-consolidation}

\begin{theorem}[Relative Mealy dynamic proarrow logic]
\label{thm:relative-mealy-dynamic-proarrow-logic}
The constructions of this chapter form the following hierarchy.

\begin{enumerate}
\item Internal categories in \(\mathcal E\) are monads in
\(\mathbf{Span}(\mathcal E)\).

\item The Eilenberg--Moore algebras of the induced monad
\[
T_{\mathbb A}:\mathcal E/A_0\to\mathcal E/A_0
\]
are right internal \(\mathbb A\)-actions. Their external category is denoted
\[
\mathbf{Act}(\mathbb A).
\]

\item A Mealy morphism
\[
\Phi:\mathbb A\rightsquigarrow\mathbb B
\]
induces a stateful restriction functor
\[
\Phi^\ast:\mathbf{Act}(\mathbb B)\to\mathbf{Act}(\mathbb A).
\]

\item For every right \(\mathbb B\)-action \(Y\), the induced right
\(\mathbb A\)-action \(\Phi^\ast Y\) has an action category
\[
\mathsf{Prog}_{\Phi,Y}
=
\int_{\mathbb A}\Phi^\ast Y.
\]

\item The relative state object is
\[
S_{\Phi,Y}=P\times_{B_0}Y.
\]
The primitive execution span
\[
\mathbf m_{\Phi,Y}:S_{\Phi,Y}\nrightarrow S_{\Phi,Y}
\]
is a monoid object in
\[
\operatorname{EndSpan}_{\mathcal E}(S_{\Phi,Y}).
\]

\item There are structural functors
\[
I_{\Phi,Y}:\mathsf{Prog}_{\Phi,Y}\to\mathbb A^{\mathrm{op}}
\]
and
\[
\Lambda_{\Phi,Y}:\mathsf{Prog}_{\Phi,Y}\to\int_{\mathbb B}Y.
\]

\item If products of internal categories exist, these assemble into a combined
labelling functor
\[
L_{\Phi,Y}:
\mathsf{Prog}_{\Phi,Y}
\to
\mathbb A^{\mathrm{op}}\times\int_{\mathbb B}Y.
\]

\item The assignment
\[
Y\longmapsto\mathsf{Prog}_{\Phi,Y}
\]
is functorial in the downstream \(\mathbb B\)-action \(Y\).

\item A Mealy transformation
\[
\theta:\Phi\Rightarrow\Psi
\]
induces, for every \(Y\), a labelled internal functor
\[
\Theta_Y:\mathsf{Prog}_{\Phi,Y}\to\mathsf{Prog}_{\Psi,Y}
\]
over both the input projection and downstream comparison.

\item Mealy morphisms
\[
\mathbb A\rightsquigarrow\mathbb B
\]
are equivalently terminal labelled input-discrete categories
\[
\mathbb P\to\mathbb A^{\mathrm{op}},
\qquad
\mathbb P\to\mathbb B^{\mathrm{op}}.
\]
This equivalence is an equivalence of data and is not a recovery of ordinary
internal natural transformations.

\item For the terminal right \(\mathbb B\)-action
\[
\mathbf 1_{\mathbb B},
\]
one obtains the trace program category
\[
\mathsf{Prog}_\Phi
=
\mathsf{Prog}_{\Phi,\mathbf 1_{\mathbb B}},
\]
and the downstream comparison functor becomes the output labelling functor
\[
\Omega_\Phi:\mathsf{Prog}_\Phi\to\mathbb B^{\mathrm{op}}.
\]

\item Representable Mealy morphisms induced by internal functors are the
stateless cases. In the terminal trace case, their program category is
\[
\mathbb A^{\mathrm{op}}.
\]

\item If \(\mathcal E\) is regular, every relative execution span acts on
subobject predicates by the diamond
\[
\langle R\rangle\varphi=\exists_\ell(r^\ast\varphi).
\]

\item If \(\mathcal E\) is coherent, finite joins of relative subprograms give
finite nondeterministic choice and satisfy
\[
\langle U\vee V\rangle\varphi
=
\langle U\rangle\varphi
\vee
\langle V\rangle\varphi.
\]

\item If \(\mathcal E\) is first-order Heyting, every relative execution span
also acts by the box
\[
[R]\varphi=\forall_\ell(r^\ast\varphi),
\]
and tests satisfy
\[
[\theta?]\varphi=\theta\Rightarrow\varphi.
\]

\item If \(\mathcal E\) has stable countable coproducts, every relative
endoprogram
\[
U:S_{\Phi,Y}\nrightarrow S_{\Phi,Y}
\]
has a finite-path star
\[
U^\star_{\mathrm{path}}
=
\coprod_{n\ge0}U^n
\]
satisfying the recursive equation
\[
U^\star_{\mathrm{path}}
\cong
1_{S_{\Phi,Y}}\oplus U;U^\star_{\mathrm{path}}.
\]

\item This path star is the free monoid on \(U\) in
\[
\operatorname{EndSpan}_{\mathcal E}(S_{\Phi,Y}).
\]

\item If \(\mathcal E\) is \(\omega\)-iterative first-order Heyting, witnessful
path semantics agrees with fibrewise fixed-point semantics:
\[
\langle U^\star_{\mathrm{path}}\rangle\varphi
=
\mu X.\bigl(\varphi\vee\langle U\rangle X\bigr),
\]
and
\[
[U^\star_{\mathrm{path}}]\varphi
=
\nu X.\bigl(\varphi\wedge[U]X\bigr).
\]

\item The corresponding induction and coinduction principles are:
\[
\varphi\vee\langle U\rangle\psi\le\psi
\quad\Longrightarrow\quad
\mu X.\bigl(\varphi\vee\langle U\rangle X\bigr)\le\psi,
\]
and
\[
\psi\le\varphi\wedge[U]\psi
\quad\Longrightarrow\quad
\psi\le\nu X.\bigl(\varphi\wedge[U]X\bigr).
\]

\item If
\[
U\hookrightarrow \mathbf m_{\Phi,Y}
\]
is a relative subprogram, the folded trace map
\[
\operatorname{fold}_{U,Y}:
U^\star_{\mathrm{path}}\to\mathbf m_{\Phi,Y}
\]
is the unique monoid morphism extending the inclusion of \(U\). It folds a
finite path of primitive relative executions into a single composite arrow of
\[
\mathsf{Prog}_{\Phi,Y}.
\]
\end{enumerate}
\end{theorem}

\begin{proof}
Part (1) is the span-monoid presentation of internal categories. Parts (2) and
(3) are Proposition~\ref{prop:actions-are-em-algebras} and
Theorem~\ref{thm:mealy-restriction-functor}. Parts (4) and (5) are
Theorem~\ref{thm:relative-program-category}. Parts (6) and (7) are
Theorem~\ref{thm:input-and-downstream-functors} and
Corollary~\ref{cor:combined-relative-labelling}. Part (8) is
Theorem~\ref{thm:program-categories-functorial-in-Y}. Part (9) is
Proposition~\ref{prop:relative-functoriality-mealy-transformations}. Part
(10) is Theorem~\ref{thm:mealy-as-labelled-input-discrete-relative} together
with Remark~\ref{rem:labelled-input-discrete-not-ordinary-2-theory}. Part
(11) is Corollary~\ref{cor:terminal-trace-program-category} and
Corollary~\ref{cor:terminal-output-labelling}. Part (12) is
Corollary~\ref{cor:terminal-stateless-case}. Part (13) is
Proposition~\ref{prop:diamond-functoriality-relative} and
Theorem~\ref{thm:universal-diamond-relative}. Part (14) is
Proposition~\ref{prop:relative-choice-law}. Part (15) is
Proposition~\ref{prop:box-functoriality-relative},
Theorem~\ref{thm:universal-box-relative}, and
Proposition~\ref{prop:relative-box-test-law}. Part (16) is
Proposition~\ref{prop:relative-path-star-equation}. Part (17) is
Theorem~\ref{thm:relative-path-star-free-monoid}. Part (18) is
Theorem~\ref{thm:relative-path-fixed-point-comparison}. Part (19) is
Corollary~\ref{cor:relative-diamond-induction} and
Corollary~\ref{cor:relative-box-coinduction}. Part (20) is
Definition~\ref{def:relative-folded-trace-map} and
Proposition~\ref{prop:relative-fold-maps-span-maps}.
\end{proof}

\begin{remark}[Hierarchy of structure]
\label{rem:relative-hierarchy-of-structure}
The constructions appear at increasing levels of categorical strength:
\[
\begin{array}{c|c|c}
\text{level} & \text{categorical structure} & \text{construction} \\
\hline
0 & \text{pullbacks} &
\Phi^\ast:\mathbf{Act}(\mathbb B)\to\mathbf{Act}(\mathbb A) \\
1 & \text{pullbacks} &
\mathsf{Prog}_{\Phi,Y}=\int_{\mathbb A}\Phi^\ast Y \\
2 & \text{regular images} &
\diamond\text{-modalities on }\mathrm{Sub}(S_{\Phi,Y}) \\
3 & \text{coherent finite joins} &
\text{finite choice of relative subprograms} \\
4 & \text{right adjoints to pullback} &
\Box\text{-modalities and implication tests} \\
5 & \text{stable countable coproducts} &
U^\star_{\mathrm{path}} \\
6 & \text{complete fibres and continuity} &
\mu,\nu\text{ fixed-point comparison}
\end{array}
\]
Thus the Mealy morphism supplies the stateful operational restriction of
behaviours, while the ambient category supplies the progressively stronger
logical operations.
\end{remark}

\begin{remark}[Final conceptual summary]
\label{rem:relative-final-summary}
A Mealy morphism
\[
\Phi:\mathbb A\rightsquigarrow\mathbb B
\]
is a stateful interpretation of \(\mathbb A\)-inputs as \(\mathbb B\)-outputs.
Dynamic logic requires a model in which outputs can act. Therefore, for each
right \(\mathbb B\)-action \(Y\), the Mealy morphism induces a right
\(\mathbb A\)-action
\[
\Phi^\ast Y.
\]
The relative program category
\[
\mathsf{Prog}_{\Phi,Y}
=
\int_{\mathbb A}\Phi^\ast Y
\]
is the execution category of the Mealy adapter coupled to the downstream
\(\mathbb B\)-system \(Y\). Dynamic proarrow logic is the modal and fixed-point
logic of spans on the coupled state object
\[
P\times_{B_0}Y.
\]
The terminal action gives the trace fragment in which outputs are observed as
labels rather than executed in a nontrivial downstream system.
\end{remark}